\documentclass[12pt, a4paper]{article}

% ── Packages ──────────────────────────────────────────────────────────────────
\usepackage[utf8]{inputenc}
\usepackage[T1]{fontenc}
\usepackage{lmodern}
\usepackage{amsmath, amssymb, amsthm}
\usepackage{mathtools}          % for \coloneqq, \DeclarePairedDelimiter, etc.
\usepackage{booktabs}           % nicer tables
\usepackage{array}              % extended column types (>{\raggedright} etc.)
\usepackage{enumitem}           % custom list formatting
\usepackage{hyperref}           % clickable cross-references
\usepackage[margin=2.5cm]{geometry}
\usepackage{parskip}            % blank lines between paragraphs, no indent
\usepackage{xcolor}
\usepackage{tcolorbox}          % coloured boxes for "intuitive explanation" etc.
\usepackage{titlesec}           % section formatting
\usepackage{cancel}             % for cancel notation in proofs

% ── Colours ───────────────────────────────────────────────────────────────────
\definecolor{intuitionbg}{RGB}{240,248,255}   % pale blue
\definecolor{simulationbg}{RGB}{255,250,240}  % pale orange
\definecolor{keybox}{RGB}{245,245,220}        % pale yellow
\definecolor{historybg}{RGB}{245,240,250}     % pale purple

% ── Environments ──────────────────────────────────────────────────────────────
\newtcolorbox{intuition}{
  colback=intuitionbg, colframe=blue!40!black,
  title={\textbf{Intuitive Explanation}}, fonttitle=\bfseries
}
\newtcolorbox{simulation}{
  colback=simulationbg, colframe=orange!60!black,
  title={\textbf{Simulation}}, fonttitle=\bfseries
}
\newtcolorbox{keytakeaway}{
  colback=keybox, colframe=olive!60!black,
  title={\textbf{Key Takeaway}}, fonttitle=\bfseries
}
\newtcolorbox{plainlanguage}{
  colback=white, colframe=gray!50,
  left=6pt, right=6pt, top=4pt, bottom=4pt
}
\newtcolorbox{historynote}{
  colback=historybg, colframe=purple!40!black,
  title={\textbf{Historical Note}}, fonttitle=\bfseries
}

\newtheorem{theorem}{Theorem}
\newtheorem{definition}{Definition}
\newtheorem{proposition}{Proposition}

% ── Notation shorthand ────────────────────────────────────────────────────────
\newcommand{\E}{\mathbb{E}}            % expected value
\newcommand{\Prob}{\mathbb{P}}         % probability
\newcommand{\Var}{\operatorname{Var}}
\newcommand{\Cov}{\operatorname{Cov}}
\newcommand{\Cor}{\operatorname{Cor}}
\newcommand{\SD}{\operatorname{SD}}
\newcommand{\given}{\,\vert\,}         % conditional bar with spacing
\newcommand{\LR}{\Lambda}              % likelihood ratio
\newcommand{\R}{\mathbb{R}}


% ══════════════════════════════════════════════════════════════════════════════
\title{\textbf{Presentation Plan: From Random Walks to Sequential Testing}\\[6pt]
       \large Version~4 --- Revised After Reading All Papers}
\author{}
\date{June 2025}

\begin{document}
\maketitle
\tableofcontents
\newpage

% ══════════════════════════════════════════════════════════════════════════════
\section*{About This Document}
\addcontentsline{toc}{section}{About This Document}

\textbf{Target audience.}
People comfortable with algebra and basic probability (e.g.\ a 12th-grade
math background).  No calculus is required to follow the main ideas, though
some calculus notation (integrals) appears in later acts for completeness ---
always with plain-language translations.

\textbf{Format.}
Fourteen interactive acts (Acts~0--13) --- deliberately small steps.  Each act
has a simulation, an intuitive explanation, and a mathematical section.  The
math sections always define every symbol before using it.  Every formula is
accompanied by a plain-language translation immediately below it.

\textbf{Design philosophy.}
Better to over-explain than under-explain.  When in doubt, slow down.

\textbf{What changed from v3.}
After reading all seven reference papers in detail, several significant
corrections were made.  The most important:
\begin{itemize}[nosep]
  \item Eppo (Schmit \& Miller, 2024) does \textbf{not} use the mSPRT
        framework of Johari et al.\ (2017).  They use \textbf{confidence
        sequences} from Howard et al.\ (2021), specifically the Normal
        mixture boundary.
  \item GAVI does \textbf{not} appear in the Eppo paper.  The tuning
        parameter is $\nu$ (related to expected sample size), not $\tau$.
  \item Robbins (1970) is the bridge between Wald's SPRT and modern
        methods --- he introduced the mixture approach and confidence
        sequences.  He now receives his own subsection.
  \item CUPED details are corrected to match Deng et al.\ (2013) exactly.
  \item Wald (1945) details are enriched with his actual ASN formulas
        and savings results (47--63\%).
  \item The narrative arc is corrected: SPRT $\to$ Robbins' mixtures $\to$
        Johari's mSPRT $\to$ Howard's confidence sequences $\to$ Eppo's
        implementation.
\end{itemize}


% ══════════════════════════════════════════════════════════════════════════════
% ACT 0
% ══════════════════════════════════════════════════════════════════════════════
\newpage
\section{Act~0 --- Why Should You Care?  The Peeking Problem}

\subsection{The Motivation}

\begin{intuition}
You work at a company running an A/B test.  Half your users see the old
website, half see a new design.  Every day, you check the dashboard:
``Is the new design better?''  After a week, the dashboard shows $p=0.03$ ---
that looks significant!  You call it: the new design wins.

But here is the dirty secret of traditional statistics: \textbf{the more often
you check, the more likely you are to see a false positive.}  A test designed
to have a 5\% false positive rate can easily reach 20--30\% if you peek at it
repeatedly.  One in four or five ``winning'' experiments might actually be noise.

This is not a theoretical concern.  It is the single most common statistical
mistake in online experimentation.
\end{intuition}

\subsection{The Simulation}

\begin{simulation}
Two groups of observations (control and treatment) are generated from
\emph{identical} distributions --- there is no real effect.  A standard
$t$-test is computed after every new observation arrives.

A counter shows: ``Number of times $p<0.05$ so far.''  The user watches the
$p$-value bounce around.  Every time it dips below $0.05$, it lights up red ---
a false positive.

\textbf{Users can control:}
\begin{itemize}[nosep]
  \item How many total observations to simulate
  \item How many independent experiments to run in parallel
\end{itemize}

A bar chart accumulates: ``Fraction of experiments that showed $p<0.05$ at
\emph{some} point during peeking.''  For 1,000 simulated experiments with
regular peeking, this fraction climbs well above 5\% --- often to 20--30\%.

A second mode checks only once, at the very end.  That fraction stays near 5\%.
\end{simulation}

\subsection{Intuitive Explanation}

\begin{intuition}
Traditional statistical tests promise: ``If there is no real effect, I will
falsely cry wolf at most 5\% of the time.''  But this promise comes with fine
print: \textbf{you may only look at the result once, at a pre-determined time.}

Every time you peek, you get another chance for randomness to fool you.  It is
like rolling a die repeatedly --- the more rolls, the more likely you are to
get a six eventually, even though any single roll has only a $1/6$ chance.

The rest of this presentation builds a test that does not have this problem:
a test where you can peek as often as you like and the false positive guarantee
still holds.

That promise has a name: \textbf{anytime-valid inference}.  Getting there
requires some beautiful mathematics, which we will build up one step at a
time.
\end{intuition}

\begin{keytakeaway}
\textbf{The peeking problem:} Checking a traditional test repeatedly inflates
the false positive rate far beyond its nominal level.  We need a fundamentally
different kind of test --- one that remains valid no matter when or how often
we look.
\end{keytakeaway}


% ══════════════════════════════════════════════════════════════════════════════
% ACT 1
% ══════════════════════════════════════════════════════════════════════════════
\newpage
\section{Act~1 --- The Random Walk}

\subsection{The Simulation}

\begin{simulation}
A vertical number line appears on screen.  A dot starts at position~0.  Each
step, a coin is flipped.  Heads: the dot moves up one step.  Tails: down one
step.  The path is drawn as a graph (position vs.\ step number).

\textbf{User controls:} speed, number of steps, probability of going up
(initially fixed at $50/50$; later a slider), number of simultaneous paths.
\end{simulation}

\subsection{Intuitive Explanation}

\begin{intuition}
Imagine you are standing in the middle of a long hallway.  Flip a coin.
Heads --- step forward.  Tails --- step back.  Repeat hundreds of times.

Where will you end up?  We don't know exactly, but we can say something
precise about the \emph{range} of likely outcomes.  Most of the time you won't
wander far.  Occasionally, a long streak of heads or tails carries you far in
one direction.  The longer you walk, the wider the range of possible positions
--- but the centre of that range stays at zero.

When we run many people simultaneously, they fan out over time --- like ink
dropped into water.  This spreading has a precise mathematical shape.
\end{intuition}

\subsection{Mathematical Formulation}

\subsubsection*{Setting up the notation}

Let $X_i$ represent the outcome of the $i$th coin flip:
\[
  X_i =
  \begin{cases}
    +1 & \text{if heads (step forward),} \\
    -1 & \text{if tails (step backward).}
  \end{cases}
\]
The probability of heads is written $p$.  For a fair coin, $p = 0.5$.

\subsubsection*{Position after $n$ steps}

\[
  \boxed{S_n = \sum_{i=1}^{n} X_i = X_1 + X_2 + \cdots + X_n}
\]

\begin{plainlanguage}
Read as: ``$S$ sub $n$ equals the sum of all $X_i$, starting from $i=1$ up to
$i=n$.''
\end{plainlanguage}

\textbf{Worked example.}  First 5 flips: H, T, H, H, T, so
$X_1\!=\!+1,\; X_2\!=\!-1,\; X_3\!=\!+1,\; X_4\!=\!+1,\; X_5\!=\!-1$.
\[
  S_1=1,\quad S_2=0,\quad S_3=1,\quad S_4=2,\quad S_5=1.
\]

\subsubsection*{The average outcome (expected value)}

When $p = 0.5$, each flip has expected value zero:
\[
  \E[X_i]
    = \underbrace{0.5}_{P(\text{heads})} \times (+1)
    + \underbrace{0.5}_{P(\text{tails})} \times (-1)
    = 0.
\]
Because each step averages to zero:
\[
  \boxed{\E[S_n] = \E[X_1] + \E[X_2] + \cdots + \E[X_n] = 0}
\]

\begin{plainlanguage}
On average, across many trials, you end up back where you started.
\end{plainlanguage}

\subsubsection*{How spread out the outcomes are --- variance \& standard deviation}

\textbf{Variance} measures the average squared distance from the mean.
For a single fair coin flip:
\[
  \Var(X_i) = \E\bigl[(X_i - 0)^2\bigr]
  = 0.5 \times (+1)^2 + 0.5 \times (-1)^2
  = 1.
\]
Because flips are independent, the variance of the sum is the sum of
the variances:
\[
  \Var(S_n) = \Var(X_1) + \Var(X_2) + \cdots + \Var(X_n)
  = \underbrace{1 + 1 + \cdots + 1}_{n\text{ terms}} = n.
\]
The \textbf{standard deviation} is the square root of the variance ---
same units as the position:
\[
  \boxed{\SD(S_n) = \sqrt{\Var(S_n)} = \sqrt{n}}
\]

\begin{plainlanguage}
After 100 steps, the typical distance from zero is $\sqrt{100}=10$.
After 10,000 steps, it is $\sqrt{10{,}000}=100$.  The spread grows with the
\emph{square root} of the number of steps, not the number itself.
\end{plainlanguage}

The simulation shows the $\pm\sqrt{n}$ envelope as two curved lines.

\subsubsection*{What if the coin is not fair?}

When $p \neq 0.5$, each step has a non-zero average:
\[
  \E[X_i] = p \times (+1) + (1-p) \times (-1) = 2p - 1.
\]
For $p = 0.6$: $\E[X_i] = 0.2$.  After $n$ steps:
\[
  \E[S_n] = n(2p-1).
\]
This is a systematic \emph{drift} --- the random walk now has a trend.

\begin{keytakeaway}
\textbf{Key concepts:} stochastic process, expected value, variance, standard
deviation, the $\sqrt{n}$ growth of randomness.
\end{keytakeaway}


% ══════════════════════════════════════════════════════════════════════════════
% ACT 2
% ══════════════════════════════════════════════════════════════════════════════
\newpage
\section{Act~2 --- The Gambling Game and the Martingale}

\subsection{The Simulation}

\begin{simulation}
The same random walk, relabelled.  The vertical axis now shows
``\texteuro{} profit or loss.''  A bold horizontal line marks zero ---
break-even.

\textbf{Probability slider:} $p > 0.5$ gives the player an edge;
$p < 0.5$ gives the house an edge.

\textbf{Doubling strategy mode:} Runs 10,000~gamblers using the martingale
betting strategy.  A histogram shows the distribution of final wealth --- many
small winners, a few catastrophic losers.  A counter shows average profit
converging to zero.

\textbf{Build your own stopping rule:} Users define a rule (``Stop when ahead
by \texteuro$K$'').  The simulation runs 10,000~gamblers.  Average winnings at
stopping converge to zero.
\end{simulation}

\subsection{Intuitive Explanation}

\begin{intuition}
Win \texteuro1 for heads, lose \texteuro1 for tails.  Cumulative winnings
trace the same random walk from Act~1.

The gambler's temptation: ``I'll keep playing until I'm ahead, then quit.''
Sounds foolproof --- surely a run of luck will eventually put you in the green?

Mathematically, a fair random walk does eventually return to positive territory.
But the time it takes can be astronomically long, and while you wait you might
lose everything.

\textbf{Central insight:} There is no strategy that improves your expected
outcome in a fair game.  This is a proven mathematical theorem.  Try any
stopping rule in the simulation --- the average winnings always converge to
zero.
\end{intuition}

\begin{intuition}
\textbf{The doubling strategy (Martingale betting), exposed.}

Lose, double the bet, repeat until you win back your losses plus \texteuro1.
It works --- most of the time.

But a long losing streak forces bets of
\texteuro1, 2, 4, 8, 16, 32, 64, \ldots\
After just 10 consecutive losses, you must bet \texteuro1,024 to win back
\texteuro1.  The histogram makes it vivid: a tall bar of people who won
\texteuro1, and a barely visible bar at the far left of people who lost
thousands.  Multiply heights by positions and sum: the average is zero.
\end{intuition}

\subsection{Mathematical Formulation}

\begin{definition}[Martingale]
A sequence $M_0, M_1, M_2, \ldots$ is a \textbf{martingale} if, at every
step, the best prediction of the next value is the current value:
\[
  \boxed{\E\bigl[M_n \given M_0, M_1, \ldots, M_{n-1}\bigr] = M_{n-1}}
\]
\end{definition}

\textbf{Reading the notation carefully:}
\begin{itemize}[nosep]
  \item $\E[\cdot]$ = ``the expected value of\ldots'' (long-run average).
  \item The vertical bar ``$\given$'' = ``given that we know.''
        Everything after the bar is information we have.
  \item $M_0, M_1, \ldots, M_{n-1}$ = the full history up to step $n-1$.
  \item Full sentence: ``The expected value of $M_n$, given that we know
        everything up to step $n-1$, equals $M_{n-1}$.''
\end{itemize}

\begin{plainlanguage}
Knowing the entire history, your best guess for the next value is just the
current value.  The past gives no useful information about the future direction.
\end{plainlanguage}

\subsubsection*{Verifying: the coin-flip game is a martingale}

Let $M_n = S_n$ (cumulative winnings).  Then $M_n = M_{n-1} + X_n$.
\begin{align*}
  \E[M_n \given M_0,\ldots,M_{n-1}]
  &= \E[M_{n-1} + X_n \given M_0,\ldots,M_{n-1}] \\
  &= M_{n-1} + \E[X_n \given M_0,\ldots,M_{n-1}]
     & &\text{($M_{n-1}$ is known, passes through $\E$)} \\
  &= M_{n-1} + 0
     & &\text{(next flip is independent; $\E[X_n]=0$)} \\
  &= M_{n-1}. \qquad\checkmark
\end{align*}

\subsubsection*{Doob's Optional Stopping Theorem (1953)}

\begin{theorem}[Doob's Optional Stopping Theorem]
Let $\{M_n\}$ be a martingale and $\tau$ a \textbf{stopping time} --- a rule
for when to quit that uses only information accumulated so far (no looking into
the future).  Then, under reasonable conditions:
\[
  \boxed{\E[M_\tau] = \E[M_0]}
\]
\end{theorem}

\begin{plainlanguage}
For a gambler starting with $M_0 = 0$: expected winnings at stopping $= 0$,
always, for any strategy.  \textbf{You cannot beat a fair game by choosing
when to quit.}
\end{plainlanguage}

\subsubsection*{What is a stopping time?}

$\tau$ is a decision rule of the form: ``Stop after step $n$ if [condition
based only on what happened so far].''

\begin{center}
\begin{tabular}{@{} l c @{}}
\toprule
\textbf{Rule} & \textbf{Valid?} \\
\midrule
Stop after exactly 100 flips & \checkmark \\
Stop the first time I am ahead by \texteuro10 & \checkmark \\
Stop when I have lost \texteuro50 & \checkmark \\
Stop after 1000 flips \emph{or} when ahead by \texteuro5, whichever first & \checkmark \\
Stop right before the next loss & $\times$ (requires the future!) \\
\bottomrule
\end{tabular}
\end{center}

\subsubsection*{The ``reasonable conditions'' for Doob's theorem}

Doob's theorem requires at least one of:
\begin{enumerate}[nosep]
  \item $\tau \leq N$ for some fixed $N$ (bounded stopping time), \textbf{or}
  \item $|M_n| \leq C$ for some fixed $C$ (bounded martingale), \textbf{or}
  \item $\E[\tau] < \infty$ \emph{and} certain integrability
        conditions (finite expected stopping time).
\end{enumerate}

When \emph{none} hold --- as with the doubling strategy --- the guarantee
breaks down.

\subsubsection*{Why the doubling strategy fails}

The doubling strategy guarantees eventual profit \emph{only if} you have
infinite money and infinite time.  In practice:
\begin{itemize}[nosep]
  \item Finite budget $\Rightarrow$ real probability of ruin before recovery.
  \item Expected time to recovery is infinite: $\E[\tau] = \infty$.
  \item Doob's theorem gives $\E[M_\tau] \leq 0$, not $= 0$.
\end{itemize}

\subsubsection*{Why this matters for what is coming}

The martingale concept reappears in every subsequent act.  The likelihood
ratio (Act~4) is a martingale.  The mixture likelihood ratio (Act~7) is a
martingale.  The reason sequential tests control false positives --- even
with peeking --- is that they exploit the same structure: \textbf{a fair game
cannot be beaten by choosing when to stop.}

\begin{keytakeaway}
\textbf{Key concepts:} martingale, conditional expectation, stopping time,
Doob's Optional Stopping Theorem, the ``reasonable conditions.''
\end{keytakeaway}


% ══════════════════════════════════════════════════════════════════════════════
% ACT 3
% ══════════════════════════════════════════════════════════════════════════════
\newpage
\section{Act~3 --- Hypothesis Testing: Telling Two Coins Apart}

\subsection{The Simulation}

\begin{simulation}
The user is presented with a coin and must decide: fair or biased?

Two panels run side by side:
\begin{itemize}[nosep]
  \item \textbf{Left:} Coin is actually fair ($p=0.5$).
  \item \textbf{Right:} Coin is actually biased ($p=0.5+\delta$).
\end{itemize}
After each flip, a running tally shows heads/tails and current proportion.
For small $\delta$, the sequences look nearly identical for many flips.

\textbf{User controls:} true bias $\delta$; number of flips before deciding.
\end{simulation}

\subsection{Intuitive Explanation}

\begin{intuition}
You suspect a coin might be unfair.  Flip it many times and look at the
results.  10~flips with 6~heads?  Could be luck.  1,000~flips with 600~heads?
Strong evidence.

Statisticians frame this as two competing stories:
\begin{itemize}[nosep]
  \item \textbf{Null hypothesis $H_0$:} ``The coin is fair.  Nothing is happening.''
  \item \textbf{Alternative hypothesis $H_1$:} ``The coin is biased.  Something real is going on.''
\end{itemize}

Two types of mistakes:
\begin{enumerate}[nosep]
  \item \textbf{False positive} (Type~I): conclude biased when it is fair.
  \item \textbf{False negative} (Type~II): conclude fair when it is biased.
\end{enumerate}
Traditionally: $\alpha = 0.05$ (max false positive rate),
$\beta = 0.20$ (max false negative rate).
\textbf{Power} $= 1 - \beta = 0.80$.
\end{intuition}

\subsection{Mathematical Formulation}

\subsubsection*{The likelihood}

Suppose $n$ flips produce $k$ heads.

Under $H_0$ (fair coin, $p = 0.5$):
\[
  \mathcal{L}_0
  = (0.5)^k \times (0.5)^{n-k}
  = (0.5)^n.
\]

Under $H_1$ (biased coin, $p = 0.5 + \delta$):
\[
  \mathcal{L}_1
  = (0.5 + \delta)^k \times (0.5 - \delta)^{n-k}.
\]

\subsubsection*{The likelihood ratio}

\[
  \boxed{\LR_n
  = \frac{\mathcal{L}_1}{\mathcal{L}_0}
  = \frac{(0.5+\delta)^k \;(0.5-\delta)^{n-k}}{(0.5)^n}}
\]

\begin{plainlanguage}
``How many times more likely is our observed data under the biased-coin story
compared to the fair-coin story?''
\end{plainlanguage}

\textbf{Worked example.}
$\delta = 0.1$ (biased coin has $p=0.6$); $n=10$ flips, $k=7$ heads.
\begin{align*}
  \mathcal{L}_1 &= (0.6)^7 \times (0.4)^3 = 0.0280 \times 0.064 \approx 0.001792, \\
  \mathcal{L}_0 &= (0.5)^{10} \approx 0.000977, \\
  \LR_{10} &= \frac{0.001792}{0.000977} \approx 1.83.
\end{align*}
The data is $\approx\!1.83$ times more likely if the coin is biased.
Evidence, but not overwhelming.

\subsubsection*{Interpreting $\LR_n$}

\begin{center}
\begin{tabular}{@{} r l @{}}
\toprule
$\LR_n = 1$   & Data equally consistent with both hypotheses \\
$\LR_n = 10$  & 10$\times$ more likely under $H_1$ --- moderate evidence \\
$\LR_n = 100$ & 100$\times$ more likely under $H_1$ --- strong evidence \\
$\LR_n = 0.1$ & 10$\times$ more likely under $H_0$ --- evidence for fairness \\
\bottomrule
\end{tabular}
\end{center}

\subsubsection*{Incremental updating}

$\LR_n$ can be computed one flip at a time:
\[
  \boxed{\LR_n = \LR_{n-1} \times \frac{f_1(x_n)}{f_0(x_n)}}
\]
where $f_j(x_n)$ is the probability of the $n$th flip's outcome under $H_j$.
\begin{itemize}[nosep]
  \item Heads: multiply by $\dfrac{0.5+\delta}{0.5}$.
  \item Tails: multiply by $\dfrac{0.5-\delta}{0.5}$.
\end{itemize}
This incremental structure is what makes \emph{sequential} testing possible.

\subsubsection*{Starting value}

Before any data:
\[
  \LR_0 = 1.
\]
Neither hypothesis is favoured.  This will be important for Ville's inequality
(which requires the starting value to equal~1).

\begin{keytakeaway}
\textbf{Key concepts:} null hypothesis, alternative hypothesis, likelihood,
likelihood ratio, false positive, false negative, power, incremental updating.
\end{keytakeaway}


% ══════════════════════════════════════════════════════════════════════════════
% ACT 4
% ══════════════════════════════════════════════════════════════════════════════
\newpage
\section{Act~4 --- The Likelihood Ratio Is a Martingale}

\subsection{The Simulation}

\begin{simulation}
The likelihood ratio $\LR_n$ is now plotted as a path over time --- like the
random walk from Act~1.

\textbf{Left panel:} 100 paths of $\LR_n$ when the coin is truly \emph{fair}
($H_0$ true).  Paths wander with no systematic trend.

\textbf{Right panel:} 100 paths when the coin is truly \emph{biased}
($H_1$ true).  Paths drift upward.
\end{simulation}

\subsection{Intuitive Explanation}

\begin{intuition}
Under $H_0$, the likelihood ratio wanders up and down but never develops a
consistent trend.  On average, it stays at~1.

This \emph{is} a martingale --- exactly like the gambler's cumulative winnings.
Everything from Act~2 applies: Doob's theorem says you cannot systematically
make $\LR_n$ large by choosing clever peeking times.
\end{intuition}

\subsection{Mathematical Formulation}

\subsubsection*{Proof that $\LR_n$ is a martingale under $H_0$}

We must show:
$\E[\LR_n \given \LR_0, \ldots, \LR_{n-1}] = \LR_{n-1}$.

From Act~3, $\LR_n = \LR_{n-1} \times \dfrac{f_1(x_n)}{f_0(x_n)}$.
Therefore:
\begin{align}
  \E\bigl[\LR_n \given \text{past}\bigr]
  &= \E\!\left[\LR_{n-1} \cdot \frac{f_1(x_n)}{f_0(x_n)} \;\bigg\vert\; \text{past}\right]
     \nonumber \\
  &= \LR_{n-1} \cdot \E\!\left[\frac{f_1(x_n)}{f_0(x_n)}\right]
     & &\text{($\LR_{n-1}$ is known)}
     \label{eq:lr-mart-step1}
\end{align}

Now the key step.  Under $H_0$, $x_n$ is drawn from $f_0$.  So:
\begin{align}
  \E\!\left[\frac{f_1(x_n)}{f_0(x_n)}\right]
  &= \sum_{\text{all outcomes } x} f_0(x) \cdot \frac{f_1(x)}{f_0(x)}
     \nonumber \\[6pt]
  &= \sum_x \cancel{f_0(x)} \cdot \frac{f_1(x)}{\cancel{f_0(x)}}
     & &\text{\textbf{(the cancellation!)}}
     \nonumber \\[6pt]
  &= \sum_x f_1(x)
     \nonumber \\[4pt]
  &= 1.
     & &\text{($f_1$ is a probability distribution; probabilities sum to~1)}
     \label{eq:lr-mart-cancel}
\end{align}

Therefore:
\[
  \E\bigl[\LR_n \given \text{past}\bigr]
  = \LR_{n-1} \times 1
  = \LR_{n-1}. \qquad\checkmark
\]

\subsubsection*{The cancellation with actual numbers}

For $\delta = 0.1$ (biased coin $p=0.6$):

\medskip
\begin{center}
\begin{tabular}{@{} l c c c @{}}
\toprule
\textbf{Outcome} & $f_0(x)$ & $f_1(x)/f_0(x)$ & $f_0(x)\times\bigl[f_1(x)/f_0(x)\bigr]$ \\
\midrule
Heads & 0.5 & $0.6/0.5 = 1.2$ & $0.5 \times 1.2 = 0.6$ \\
Tails & 0.5 & $0.4/0.5 = 0.8$ & $0.5 \times 0.8 = 0.4$ \\
\midrule
\multicolumn{3}{r}{\textbf{Total:}} & $\mathbf{1.0}$ \checkmark \\
\bottomrule
\end{tabular}
\end{center}

\begin{plainlanguage}
When we compute the expected ratio under $H_0$, the $H_0$ probabilities
cancel out, and we are just left summing $H_1$ probabilities --- which must
add to~1.  This is the ``aha'' moment of the proof.
\end{plainlanguage}

\subsubsection*{Two additional properties}

\begin{enumerate}[nosep]
  \item \textbf{$\LR_n$ is non-negative:} It is a product of non-negative
        ratios (probabilities divided by probabilities), so $\LR_n \geq 0$.
  \item \textbf{$\LR_0 = 1$:} Before any data, neither hypothesis is favoured.
\end{enumerate}

So $\LR_n$ is a \textbf{non-negative martingale starting at~1}.  Remember
these properties --- they are exactly what Ville's inequality needs.

\subsubsection*{Under $H_1$, $\LR_n$ drifts upward}

When the coin really is biased:
\[
  \E\!\left[\frac{f_1(x_n)}{f_0(x_n)}\right]_{\!H_1}
  = \sum_x f_1(x) \cdot \frac{f_1(x)}{f_0(x)}
  > 1.
\]
So $\LR_n$ tends to grow --- this upward drift is the signal the test detects.

\begin{keytakeaway}
\textbf{Key concepts:} the likelihood ratio is a martingale under~$H_0$;
the cancellation that makes it work; non-negative martingale; drift under~$H_1$.
\end{keytakeaway}


% ══════════════════════════════════════════════════════════════════════════════
% ACT 5
% ══════════════════════════════════════════════════════════════════════════════
\newpage
\section{Act~5 --- Markov, Ville, and the Anytime-Valid Guarantee}

\subsection{The Simulation}

\begin{simulation}
This is the pivotal simulation of the entire presentation.

\textbf{Panel~1 --- Markov (single time):}
10,000 fair-coin paths of $\LR_n$.  A vertical line at $n=200$.
A threshold at $\LR = 1/\alpha = 20$ (for $\alpha = 0.05$).
Count: how many paths are above the threshold \emph{at step~200}?
Fraction $\leq 5\%$.

\textbf{Panel~2 --- Ville (any time):}
Same 10,000 paths.  Count: how many \emph{ever} crossed the threshold at
\emph{any} step?  Also $\leq 5\%$ --- but requires a stronger theorem.

\textbf{Panel~3 --- Comparison:}
10,000 paths of a standard $z$-score checked at every step.  Count how many
ever show significance.  Much higher than 5\% --- the peeking problem.

Summary:
\begin{itemize}[nosep]
  \item Markov: controls false positives at a single time \checkmark
  \item Ville: controls across all times simultaneously \checkmark
  \item Standard test with peeking: does NOT control $\times$
\end{itemize}
\end{simulation}

\subsection{Intuitive Explanation}

\begin{intuition}
We are about to prove the most important result in this entire presentation:
the \textbf{anytime-valid guarantee}.

Two steps:
\begin{enumerate}[nosep]
  \item \textbf{Markov's inequality} --- a simple bound at one point in time.
  \item \textbf{Ville's inequality} --- extends the bound to all points in
        time simultaneously.
\end{enumerate}
Ville's inequality is \emph{the reason} sequential testing works.
\end{intuition}

\subsection{Mathematical Formulation}

% ── Markov ─────────────────
\subsubsection*{Step 1: Markov's Inequality}

\begin{theorem}[Markov's Inequality]
For any non-negative random variable $X$ with finite $\E[X]$:
\[
  \boxed{\Prob(X \geq c) \;\leq\; \frac{\E[X]}{c}}
\]
\end{theorem}

\begin{plainlanguage}
$\Prob(\cdot)$ denotes probability --- a number between 0 and~1 measuring how
likely an event is.  So ``the probability that $X$ is at least $c$ cannot
exceed the average of $X$ divided by $c$.''
\end{plainlanguage}

\textbf{Why is this true?  An intuitive proof.}

Suppose $X \geq 0$ always.  Let $q = \Prob(X \geq c)$ be the fraction of
outcomes at or above $c$.  Those outcomes contribute \emph{at least}
$q \times c$ to the average (each is $\geq c$).  The rest contribute $\geq 0$
(since $X \geq 0$):
\[
  \E[X] \;\geq\; q \cdot c + (1-q)\cdot 0 = q \cdot c
  \qquad\Longrightarrow\qquad
  q \leq \frac{\E[X]}{c}. \qquad\checkmark
\]

\textbf{Concrete example.}
Average income in a town is \texteuro50,000.
\[
  \Prob(\text{income} \geq \text{\texteuro}500{,}000)
  \;\leq\;
  \frac{50{,}000}{500{,}000} = 0.10 = 10\%.
\]
At most 10\% of people can earn \texteuro500,000+.

\textbf{Applying Markov to the likelihood ratio.}

Under $H_0$, $\LR_n$ is a martingale with $\E[\LR_n] = \LR_0 = 1$ and
$\LR_n \geq 0$.  Markov gives:
\[
  \Prob\!\left(\LR_n \geq \frac{1}{\alpha}\right)
  \leq \frac{\E[\LR_n]}{1/\alpha}
  = 1 \cdot \alpha
  = \alpha.
\]

\begin{plainlanguage}
At any single, pre-specified step $n$, the false positive probability
is at most $\alpha$.  Reassuring --- but only for \emph{one} moment.
\end{plainlanguage}

% ── Ville ──────────────────
\subsubsection*{Step 2: Ville's Inequality (1939)}

\begin{theorem}[Ville's Inequality]
For any \textbf{non-negative martingale} $\{M_n\}_{n \geq 0}$ with
$M_0 = m_0$:
\[
  \boxed{\Prob\!\left(\sup_{n \geq 0}\, M_n \;\geq\; c\right)
         \;\leq\; \frac{m_0}{c}}
\]
where $\sup_{n\geq 0} M_n$ is the highest value $M_n$ \emph{ever} reaches.
\end{theorem}

\textbf{Why this is remarkable:}

\begin{center}
\renewcommand{\arraystretch}{1.3}
\begin{tabular}{@{} l p{9cm} @{}}
\toprule
\textbf{Markov} & At any single \emph{fixed} moment, $\Prob(M_n \geq c)$ is small. \\
\textbf{Ville}  & $\Prob(M_n$ \textbf{ever} $\geq c$, at \textbf{any} moment$)$ is \emph{also} small. \\
\bottomrule
\end{tabular}
\end{center}

Markov would allow the process to spike above~$c$ at various times, as long as
it is rarely there at any single pre-specified time.  Ville says even those
spikes are collectively rare.

\textbf{The intuition behind Ville.}
For a non-negative martingale, reaching a high value ``uses up'' the available
expected value.  Once high, it must stay high on average (martingale property),
yet it cannot go below zero (non-negativity).  These two constraints together
limit how often it can ever be high.

\subsubsection*{Applying Ville to the likelihood ratio}

Set $M_n = \LR_n$, $m_0 = \LR_0 = 1$, $c = 1/\alpha$:
\[
  \boxed{\Prob\!\left(\LR_n \text{ ever reaches } \frac{1}{\alpha}
         \text{ or higher}\right) \;\leq\; \alpha}
\]

\begin{plainlanguage}
The probability that the likelihood ratio ever crosses our rejection
threshold --- at any point, across all flips, even if you run forever ---
is at most $\alpha$.

\textbf{This is the anytime-valid guarantee.}  You can peek after every
observation.  The false positive rate, across all checks combined, stays
$\leq \alpha$.
\end{plainlanguage}

\subsubsection*{Connecting back to Act~0}

Standard test statistics (like $p$-values from a $t$-test) are \emph{not}
martingales.  Checking repeatedly gives many independent chances to be fooled.

The likelihood ratio \emph{is} a martingale under~$H_0$.  Ville's inequality
converts this into a practical guarantee: \textbf{peeking is safe.}

\begin{historynote}
Jean Ville proved this inequality in his 1939 doctoral thesis,
\textit{\'Etude critique de la notion de collectif}.  It predated the formal
development of martingale theory by Doob (1953) and was largely overlooked
for decades.  Robbins (1970) was among the first to recognise its practical
importance for sequential testing --- a connection we explore in Act~7.
\end{historynote}

\begin{keytakeaway}
\textbf{Key concepts:} Markov's inequality (with proof), Ville's inequality,
the anytime-valid guarantee, martingale + non-negative = peek-safe.
\end{keytakeaway}


% ══════════════════════════════════════════════════════════════════════════════
% ACT 6
% ══════════════════════════════════════════════════════════════════════════════
\newpage
\section{Act~6 --- SPRT: Wald's Sequential Probability Ratio Test (1945)}

\subsection{The Simulation}

\begin{simulation}
Two coins: one fair ($H_0$), one biased ($H_1$).  The user does not know
which they have.

$\LR_n$ is plotted as a path.  Two thresholds:
\begin{itemize}[nosep]
  \item \textbf{Upper} at $\LR_n = B$: cross this $\Rightarrow$ ``coin is biased.''
  \item \textbf{Lower} at $\LR_n = A$: fall below this $\Rightarrow$
        ``coin is fair.''
\end{itemize}
The path drifts toward one boundary and eventually crosses.

\textbf{User controls:} which coin is used, $\alpha$, $\beta$, $\delta$.

\textbf{Additional panel:} Distribution of stopping times across 10,000
simulated experiments, alongside the fixed-sample size required for the same
$\alpha$ and $\beta$.  The average SPRT stopping time is visibly less.
\end{simulation}

\subsection{Intuitive Explanation}

\begin{intuition}
Abraham Wald (1943, published 1945) asked: given that we can peek at the data
anytime, what is the most \emph{efficient} way to decide between two
hypotheses?

His answer: compute $\LR_n$ after each observation.  If it gets high enough,
conclude $H_1$.  Low enough, conclude $H_0$.  Otherwise, keep collecting.

Wald proved this is \textbf{optimal}: no other sequential test with the same
error guarantees needs fewer observations on average.  In fact, the SPRT
typically requires \textbf{47--63\% fewer observations} than the best
fixed-sample test with the same error rates.
\end{intuition}

\subsection{Mathematical Formulation}

\subsubsection*{Decision rule}

Choose error tolerances $\alpha$ (false positive rate) and $\beta$ (false
negative rate).  Compute two thresholds:
\[
  \boxed{B = \frac{1-\beta}{\alpha}}, \qquad
  \boxed{A = \frac{\beta}{1-\alpha}}.
\]

\textbf{Example} ($\alpha = 0.05$, $\beta = 0.20$):
\[
  B = \frac{0.80}{0.05} = 16, \qquad
  A = \frac{0.20}{0.95} \approx 0.211.
\]

\begin{center}
\begin{tabular}{@{} l l l @{}}
\toprule
\textbf{Condition} & \textbf{Decision} & \textbf{Meaning} \\
\midrule
$\LR_n \geq B = 16$        & Reject $H_0$  & ``Coin is biased'' \\
$\LR_n \leq A \approx 0.21$  & Accept $H_0$  & ``Coin is fair'' \\
$A < \LR_n < B$            & Keep flipping & ``Not enough evidence yet'' \\
\bottomrule
\end{tabular}
\end{center}

\subsubsection*{Where the thresholds come from (Wald's inequalities)}

These are \textbf{Wald's approximations}, derived from fundamental
inequalities he proved in his 1945 paper.  Let $\alpha' = \Prob_{H_0}(\text{reject } H_0)$
be the actual Type~I error and $\beta' = \Prob_{H_1}(\text{accept } H_0)$ the
actual Type~II error.  Wald showed:
\begin{align*}
  \alpha' &\leq \frac{1}{B}, \\
  \beta' &\leq A.
\end{align*}

Setting $B = (1-\beta)/\alpha$ and $A = \beta/(1-\alpha)$ guarantees
$\alpha' \leq \alpha/(1-\beta)$ and $\beta' \leq \beta/(1-\alpha)$.  When
$\alpha$ and $\beta$ are small (the typical case), these are very close to
$\alpha$ and $\beta$ themselves.

\begin{plainlanguage}
The thresholds are set so that if the likelihood ratio ever climbs to $B$ or
drops to $A$, we can be confident in our decision.  The approximation is
slightly conservative --- actual error rates are at or below the targets.
\end{plainlanguage}

\subsubsection*{The random walk view}

Equivalently, define the log-likelihood ratio $Z_n = \log \LR_n =
\sum_{i=1}^n \log \frac{f_1(x_i)}{f_0(x_i)}$, where $\log$ denotes the
natural logarithm (the inverse of the exponential function: if $e^y = x$
then $\log x = y$).  The SPRT becomes a random
walk between two barriers: stop at the first $n$ such that $Z_n \geq \log B$
or $Z_n \leq \log A$.

\subsubsection*{Expected sample size (Wald's ASN formula)}

Using what is now called \textbf{Wald's equation} --- $\E[Z_n] = \E[n]
\cdot \E[z]$, where $z = \log \frac{f_1(x)}{f_0(x)}$ is the single-observation
log-likelihood ratio --- Wald derived:
\[
  \boxed{E_0[n] \approx \frac{(1-\alpha)\log A + \alpha \log B}{\E_0[z]}}
\]
under $H_0$, and
\[
  \boxed{E_1[n] \approx \frac{\beta \log A + (1-\beta)\log B}{\E_1[z]}}
\]
under $H_1$.

\begin{plainlanguage}
The expected number of observations depends on the log-thresholds ($\log A$
and $\log B$) and how much information each observation provides (measured by
$\E[z]$, which depends on how different $H_0$ and $H_1$ are).
\end{plainlanguage}

\subsubsection*{Savings over fixed-sample tests}

Wald computed the fixed-sample size needed for the same $\alpha$ and $\beta$,
and compared it to the SPRT's expected sample size.  His results (Table~1 in
the 1945 paper):

\begin{center}
\begin{tabular}{@{} c c c @{}}
\toprule
$\alpha$ & $\beta$ & \textbf{Saving in expected observations} \\
\midrule
0.05 & 0.05 & $\approx 53\%$ \\
0.05 & 0.01 & $\approx 47\%$ \\
0.01 & 0.01 & $\approx 63\%$ \\
\bottomrule
\end{tabular}
\end{center}

\begin{plainlanguage}
The SPRT needs, on average, only about \textbf{half as many observations} as a
traditional fixed-sample test.  This was a remarkable result in 1945 and
remains one of the main motivations for sequential testing.
\end{plainlanguage}

\subsubsection*{Optimality of the SPRT}

\begin{historynote}
Wald conjectured in his 1945 paper (footnote~12, p.~157) that the SPRT is
\emph{exactly} optimal, not merely approximately so.  He proved it only for
the special case where the log-likelihood ratio jumps in fixed increments
(zero overshoot).  Three years later, Wald \& Wolfowitz
(1948) proved the general result: among \textbf{all} sequential tests
(not just SPRTs) with the same or smaller error probabilities $(\alpha, \beta)$,
the SPRT minimises the expected sample size under both $H_0$ and~$H_1$
simultaneously.
\end{historynote}

\subsubsection*{Truncation}

In practice, experiments cannot run forever.  Wald showed that truncating the
SPRT at a maximum number of observations $n_0$ barely affects error rates.
For example, with $n_0 = 2N$ (where $N$ is the fixed-sample size), the
increase in $\alpha$ and $\beta$ is negligible.

\subsubsection*{The critical weakness}

SPRT requires specifying $\delta$ in advance --- the exact alternative
hypothesis $H_1$.  If you set $\delta = 0.1$ but
the true bias is $0.03$, the likelihood ratio is computed against the wrong
alternative.  The test may take much longer, or the test may still determine
the correct answer but far less efficiently.

In A/B testing --- where the effect size is typically unknown --- this is a
serious limitation.  The next two acts address this.

\begin{keytakeaway}
\textbf{Key concepts:} SPRT decision rule, Wald's boundary approximations,
Wald's equation for expected sample size, 47--63\% saving over fixed-sample
tests, optimality (Wald \& Wolfowitz 1948), the critical dependence on a
pre-specified effect size.
\end{keytakeaway}


% ══════════════════════════════════════════════════════════════════════════════
% ACT 7
% ══════════════════════════════════════════════════════════════════════════════
\newpage
\section{Act~7 --- The Mixture Approach: From Robbins (1970) to the mSPRT}

This act has two parts.  First, the key idea: mixing over effect sizes.
Second, the application to A/B testing by Johari et al.\ (2017).

\subsection{The Simulation}

\begin{simulation}
\textbf{Panel~1 --- The problem with wrong $\delta$:}
Three simultaneous paths: SPRT with the correct $\delta$, SPRT with a
$\delta$ that is too large, and SPRT with a $\delta$ that is too small.
The correct one is fastest; the wrong ones are slow or stuck.

\textbf{Panel~2 --- The mixture idea:}
Like Act~6, but the user no longer specifies $\delta$.  A new slider
controls $\tau$ (width of a ``prior belief'' over effect sizes).
$\LR_n^{H}$ is plotted.  Only the \emph{upper} threshold $1/\alpha$ is
shown (no lower boundary).

\textbf{Panel~3 --- Comparison:}
A power curve and average stopping time chart comparing SPRT (with correct
$\delta$) to the mSPRT across a range of true effect sizes.  The mSPRT is
slightly less efficient at any single $\delta$ but much more robust overall.
\end{simulation}

\subsection{Part A: Robbins' Mixture Approach (1970)}

\subsubsection*{Intuitive Explanation}

\begin{intuition}
Wald's SPRT is like betting everything on one horse: you pick a specific
effect size $\delta$ and compute evidence based on that choice.  If the horse
wins (the true effect is near $\delta$), you are fast.  Otherwise, slow.

Herbert Robbins proposed a radical alternative in 1970: \textbf{bet on all
horses simultaneously.}  Instead of choosing one $\delta$, spread your
``evidence budget'' across many possible effect sizes using a weighting
scheme (a mixing distribution $H$).

The key insight: this average likelihood ratio is \emph{still} a non-negative
martingale under $H_0$.  Ville's inequality still applies.  The anytime-valid
guarantee is preserved.
\end{intuition}

\subsubsection*{The fundamental result}

Robbins built on Ville's inequality directly.  He observed that if $\LR_n^\delta$
is the likelihood ratio for a specific effect size $\delta$, then the
\textbf{mixture likelihood ratio}
\[
  \boxed{\LR_n^{H} = \int \LR_n^{\delta} \, dH(\delta)}
\]
is also a non-negative martingale starting at~1 under $H_0$, because:
\begin{enumerate}[nosep]
  \item Each $\LR_n^\delta$ is a non-negative martingale under $H_0$ (Act~4).
  \item A weighted average (integral) of martingales is a martingale
        (linearity of conditional expectation).
\end{enumerate}

Therefore Ville's inequality gives:
\[
  \Prob\!\left(\LR_n^{H} \text{ ever} \geq \frac{1}{\alpha}\right) \leq \alpha.
\]

\begin{plainlanguage}
Symbol by symbol:
\begin{itemize}[nosep]
  \item $\LR_n^{\delta}$ = the likelihood ratio assuming effect size $\delta$.
  \item $dH(\delta)$ = how much weight to give effect size $\delta$.
  \item $\int$ = continuous version of summation --- averages over all
        possible $\delta$.
\end{itemize}
Think: ``average likelihood ratio across an infinitely large expert panel,
weighted by credibility.''
\end{plainlanguage}

\subsubsection*{Robbins' confidence sequences}

Robbins went further.  He showed that inverting the mixture test gives a
\textbf{confidence sequence} --- a sequence of intervals $I_n$ that covers
the true parameter at \emph{all} sample sizes simultaneously:
\[
  \Prob\!\bigl(\mu \in I_n \text{ for every } n \geq 1\bigr) \geq 1 - \alpha.
\]

For the normal case with observations $x_i \sim \mathcal{N}(\mu, 1)$ and
$S_n = \sum_{i=1}^n x_i$, he derived the two-sided confidence sequence:
\[
  \boxed{I_n = \left(\frac{S_n - c_n}{n},\; \frac{S_n + c_n}{n}\right)}
\]
where
\[
  c_n = \sqrt{(n+m)\!\left(a^2 + \log\!\left(\frac{n}{m}+1\right)\right)}
\]
and $\Prob(\mu \notin I_n \text{ for some } n) \leq e^{-a^2}$.
Here $m > 0$ is a free parameter and $a^2$ controls the coverage (e.g.\
$a^2 = 6$ gives $\geq 95\%$ simultaneous coverage).

\begin{plainlanguage}
Unlike a standard confidence interval (valid at one pre-specified time),
a confidence sequence is valid at every time simultaneously.  You can peek
as often as you want.  Robbins proved this using exactly the mixture
likelihood approach --- Ville's inequality does the heavy lifting.
\end{plainlanguage}

\subsubsection*{Connection to the law of the iterated logarithm}

\begin{historynote}
Robbins showed that the \emph{tightest possible} confidence sequence boundary
grows as $\sqrt{2n \log \log n}$ --- exactly the rate of the \textbf{law of
the iterated logarithm} (LIL), which describes the furthest excursions of
a random walk.  Wider boundaries (like the $\sqrt{n \log n}$ rate from the
Normal mixing distribution) give tighter probability bounds but larger
intervals at any finite $n$.  The practical tradeoff: LIL-rate boundaries are
asymptotically optimal but wider for moderate $n$; the Normal mixture boundaries
are slightly wider asymptotically but tighter over the 2--3 orders of
magnitude of sample sizes relevant to A/B testing.
\end{historynote}


\subsection{Part B: The mSPRT for A/B Testing (Johari et al., 2017)}

\subsubsection*{Intuitive Explanation}

\begin{intuition}
Johari, Pekelis, and Walsh (2017) applied Robbins' idea to the specific
setting of A/B tests.  They chose the Normal distribution as the mixing
distribution $H$ and derived a closed-form formula for the mixture likelihood
ratio.  They called it the \textbf{mixture Sequential Probability Ratio Test
(mSPRT)}.

Their key contribution: making the theory practical.  They gave explicit
formulas, calibration guidance for $\tau$, and proved that being within
an order of magnitude of the true effect-size variance is sufficient for
near-optimal power.
\end{intuition}

\subsubsection*{Closed form for Normal observations}

Assumptions:
\begin{itemize}[nosep]
  \item Observations $x_i \sim \mathcal{N}(\mu, \sigma^2)$
        (Normal with mean $\mu$, \textbf{known} variance $\sigma^2$).
  \item Mixing distribution $H = \mathcal{N}(0, \tau^2\sigma^2)$
        (Normal centred at~0, variance proportional to $\sigma^2$).
\end{itemize}

The mixture likelihood ratio has an exact closed form (Johari et al.,
Proposition~1):
\[
  \boxed{\LR_n^{H}
  = \frac{1}{\sqrt{1 + n\tau^2}}
    \;\exp\!\left(
      \frac{n\tau^2\, \bar{x}_n^2}
           {2\,(\sigma^2/n)\,(1 + n\tau^2)}
    \right)}
\]

\textbf{Unpacking every piece:}

\begin{center}
\renewcommand{\arraystretch}{1.4}
\begin{tabular}{@{} >{\raggedright}p{3cm} p{9cm} @{}}
\toprule
\textbf{Symbol} & \textbf{Meaning} \\
\midrule
$\sigma^2$ & Variance of individual observations (how noisy the data is). \\
$n$ & Number of observations so far. \\
$\tau^2$ & Variance of the mixing distribution $H$, expressed as a multiple
  of $\sigma^2$.  Controls how wide a range of effect sizes we consider
  plausible. \\
$\bar{x}_n$ & Sample mean after $n$ observations:
  $\bar{x}_n = \frac{1}{n}\sum_{i=1}^n x_i$. \\
$1/\sqrt{1 + n\tau^2}$ & Shrinkage factor.  Slowly decreases as $n$ grows;
  prevents the statistic from growing merely because more data arrived. \\
$\exp(\cdot)$ & $e \approx 2.718$ raised to the argument.  Grows rapidly
  when $\bar{x}_n$ is far from zero. \\
\bottomrule
\end{tabular}
\end{center}

\begin{plainlanguage}
``The mSPRT score is large when the observed average is far from zero (strong
signal), there are many observations, and the noise is low relative to the
prior width.''
\end{plainlanguage}

\subsubsection*{Decision rule}

\[
  \text{Stop and reject } H_0 \text{ when:} \quad
  \boxed{\LR_n^{H} \;\geq\; \frac{1}{\alpha}}
\]

For $\alpha = 0.05$: stop when $\LR_n^{H} \geq 20$.

\textbf{No lower boundary.}
Unlike the SPRT, the mSPRT cannot definitively confirm $H_0$.  The mixture
over all effect sizes, including very small ones, means the statistic hovers
near~1 indefinitely under $H_0$.  In practice, experiments run for a fixed
maximum duration; if the threshold is never crossed, conclude: ``No significant
effect detected.''

\subsubsection*{The role of $\tau$}

\begin{center}
\begin{tabular}{@{} l l l @{}}
\toprule
$\tau$ is\ldots & \textbf{Effect} & \textbf{Best for} \\
\midrule
Small ($\sim\!0.01$) & Conservative, slow to reject & Very small effects \\
Large ($\sim\!0.5$)  & Aggressive, fast for large effects & Large effects \\
Matched to true $\delta$ & Near-optimal & Good prior knowledge \\
\bottomrule
\end{tabular}
\end{center}

Johari et al.\ (Theorem~3) proved that the optimal $\tau$ is chosen to
minimise the expected stopping time for a given true effect size.  In practice,
getting \emph{within an order of magnitude} of the true effect-size variance
suffices for near-optimal power.

\subsubsection*{Key assumption: known variance}

The mSPRT formula assumes $\sigma^2$ is known.  For continuous data this is
typically estimated; for binary data (e.g.\ click/no-click), Johari et al.\
note that the exact mSPRT results hold only approximately via the Central Limit
Theorem.  For large samples ($n \geq 100$, nearly always
true in A/B testing), the approximation is excellent.

\subsubsection*{Always-valid $p$-values and confidence intervals}

Johari et al.\ also defined \textbf{always-valid $p$-values}:
\[
  p_n = \min\!\left(1,\; \frac{1}{\LR_n^{H}}\right)
\]
and \textbf{always-valid confidence intervals} by inverting $\LR_n^H$:
\[
  C_n = \bigl\{\mu : \LR_n^{H}(\mu) < 1/\alpha \bigr\}.
\]
These enjoy the usual sequential guarantee: at any stopping time $\tau$,
the Type~I error and coverage are controlled.

\begin{keytakeaway}
\textbf{Key concepts:} Robbins' mixture idea, the mSPRT as a weighted average
of likelihood ratios, the closed-form Normal formula, the role of $\tau$,
no lower boundary, the known-variance assumption.

\textbf{What's next:}  The mSPRT was a landmark --- but it is not what Eppo
uses.  The next act introduces \textbf{confidence sequences} from
Howard et al.\ (2021), a more general framework that Eppo adopted instead.
\end{keytakeaway}


% ══════════════════════════════════════════════════════════════════════════════
% ACT 8
% ══════════════════════════════════════════════════════════════════════════════
\newpage
\section{Act~8 --- Confidence Sequences: The Modern Framework\\ (Howard et al., 2021)}

\subsection{The Simulation}

\begin{simulation}
\textbf{Panel~1 --- Confidence Sequence vs.\ Fixed CI:}
Data arrives one observation at a time.  Two bands track the estimated mean:
\begin{itemize}[nosep]
  \item A \textbf{fixed 95\% CI} (recomputed at each step) --- shrinks as
        $1/\sqrt{n}$, but invalidated by peeking.
  \item A \textbf{confidence sequence} --- wider at first, but valid at
        \emph{all} times simultaneously.
\end{itemize}
A counter shows how often the true mean escapes each band (across 10,000
runs).  For the fixed CI with peeking: $\gg 5\%$.  For the CS: $\leq 5\%$.

\textbf{Panel~2 --- The Normal mixture boundary:}
The boundary function $u(v)$ plotted against the variance $v = n\sigma^2$.
A slider for $\nu$ shows how changing $\nu$ shifts the ``sweet spot'' ---
the sample size where the boundary is tightest relative to the fixed-sample
CI.
\end{simulation}

\subsection{Intuitive Explanation}

\begin{intuition}
A confidence sequence is the natural answer to the peeking problem for
estimation (not just testing).

A standard 95\% confidence interval says: ``We are 95\% confident the true
value is in this range \emph{right now}.''  But if you check again tomorrow
with new data, the guarantee resets.  Check 100 times and you are almost
guaranteed to be wrong at least once.

A confidence sequence says: ``We are 95\% confident the true value has been
in this range \emph{at every single time we have ever checked, and at every
time we will ever check in the future}.''

This is a strictly stronger guarantee.  The price: the band is wider than a
fixed CI at any single time.  But in exchange, you are free to monitor
continuously, stop whenever you want, and report at any time.
\end{intuition}

\subsection{Mathematical Formulation}

\subsubsection*{Definition}

\begin{definition}[Confidence Sequence]
A sequence of intervals $(C_t)_{t \geq 1}$ is a $(1-\alpha)$-\textbf{confidence
sequence} (CS) for a parameter $\mu$ if:
\[
  \boxed{\Prob\!\bigl(\mu \in C_t \text{ for all } t \geq 1\bigr) \;\geq\; 1-\alpha}
\]
\end{definition}

\begin{plainlanguage}
The \textbf{entire sequence} of intervals simultaneously covers the true
parameter with probability at least $1-\alpha$.  Not just at one time ---
at every time.
\end{plainlanguage}

Compare with a standard CI, which only guarantees $\Prob(\mu \in C_n) \geq
1-\alpha$ for a single fixed~$n$.

\subsubsection*{The sub-$\psi$ framework}

Howard et al.\ build confidence sequences from \textbf{sub-$\psi$ processes},
which generalise the concept of ``observations bounded by a moment generating
function.''  The key idea:

\begin{enumerate}[nosep]
  \item Define a process $S_t$ (e.g.\ the running sum of centred observations).
  \item Verify that the process satisfies a \textbf{sub-$\psi$ condition}:
        the moment generating function of $S_t - S_{t-1}$, conditional on the
        past, is bounded by a known function $\psi$.
  \item Construct a \textbf{supermartingale} from this bound.
  \item Apply Ville's inequality to get the anytime-valid guarantee.
\end{enumerate}

For sub-Gaussian observations (which includes all bounded random variables
and all Normal random variables), the relevant condition is:
\[
  \E\!\bigl[e^{\lambda(X_t - \mu)} \given \text{past}\bigr]
  \;\leq\;
  e^{\lambda^2 \sigma^2 / 2}
  \qquad \text{for all } \lambda \in \R.
\]
This says: the tails of each observation are no heavier than those of a Normal
distribution with variance $\sigma^2$.

\begin{plainlanguage}
The sub-$\psi$ condition is a way of saying ``the data is not too wild.''
For Normal data or bounded data (like click rates), this condition is
automatically satisfied.  It allows the theory to work \emph{without}
assuming the data is exactly Normal --- making it \textbf{nonparametric}.
\end{plainlanguage}

\subsubsection*{The Normal Mixture Boundary (the key formula)}

The centrepiece of Howard et al.\ is the \textbf{Normal mixture confidence
sequence} (Proposition~5, equation~14 in their paper).  For sub-Gaussian
observations with intrinsic time (cumulative variance) $v_t = \sum_{i=1}^t
\sigma_i^2$:

\[
  \boxed{u(v) = \sqrt{(v + \nu) \cdot \log\!\frac{v + \nu}{\nu \alpha^2}}}
\]

The $(1-\alpha)$-confidence sequence is:
\[
  C_t = \left(\bar{x}_t - \frac{u(v_t)}{t},\; \bar{x}_t + \frac{u(v_t)}{t}
        \right)
\]

where $v_t = t\sigma^2$ for i.i.d.\ observations with variance $\sigma^2$.

\textbf{Unpacking the boundary:}

\begin{center}
\renewcommand{\arraystretch}{1.4}
\begin{tabular}{@{} >{\raggedright}p{2.5cm} p{9.5cm} @{}}
\toprule
\textbf{Symbol} & \textbf{Meaning} \\
\midrule
$v$ & \textbf{Intrinsic time}: the cumulative variance of observations so far.
  For i.i.d.\ data: $v_t = t\sigma^2$.  For heterogeneous data: $v_t =
  \sum_{i=1}^t \sigma_i^2$. \\
$\nu > 0$ & \textbf{Tuning parameter}: controls \emph{which sample sizes} the
  boundary is tightest for.  Specifically, the boundary is approximately
  optimised for $v \approx \nu$.  Larger $\nu$ = tighter for larger samples,
  wider for smaller samples. \\
$\alpha$ & Significance level (e.g.\ 0.05). \\
$\log(\cdot)$ & Natural logarithm. \\
$\sqrt{(v+\nu) \cdot \log\frac{v+\nu}{\nu\alpha^2}}$
  & The boundary grows roughly as $\sqrt{v \log v}$ for large $v$.  This is
  slightly wider than the $\sqrt{v}$ growth of a fixed-sample CI, but the
  extra $\sqrt{\log v}$ is the ``price of peeking.'' \\
\bottomrule
\end{tabular}
\end{center}

\begin{plainlanguage}
The boundary $u(v)$ tells you how far from zero the running sum $S_t$ can
drift before you can reject $H_0$.  The parameter $\nu$ lets you tune where
the boundary is tightest.  If you expect to collect $N$ observations total with
variance $\sigma^2$, setting $\nu \approx N\sigma^2$ makes the boundary
tightest around $n = N$.
\end{plainlanguage}

\subsubsection*{How $\nu$ works}

\begin{center}
\begin{tabular}{@{} l l l @{}}
\toprule
$\nu$ is\ldots & \textbf{Effect} & \textbf{Best for} \\
\midrule
Small & Tight for small samples, wide for large & Early stopping \\
Large & Wide for small samples, tight for large & Long experiments \\
$\approx N\sigma^2$ & Near-optimal at sample size $N$ & Known experiment length \\
\bottomrule
\end{tabular}
\end{center}

\subsubsection*{Connection to Robbins and the mSPRT}

Howard et al.\ explicitly credit Robbins (1970) as their starting point.
Their Normal mixture boundary comes from choosing a \emph{Normal mixing
distribution} over the mean parameter --- exactly Robbins' approach.  The
connection:
\begin{itemize}[nosep]
  \item Robbins (1970): introduced the mixture approach and confidence
        sequences, but only for known-variance Normal data.
  \item Johari et al.\ (2017): made it practical for A/B testing via the
        mSPRT, but still assumed known variance and parametric (Normal) data.
  \item \textbf{Howard et al.\ (2021): extended to nonparametric settings}
        via the sub-$\psi$ framework.  The boundary works for \emph{any}
        sub-Gaussian data (not just Normal), the variance can be unknown
        (estimated from data), and the coverage guarantee is exact
        (non-asymptotic).
\end{itemize}

\subsubsection*{Growth rate comparison}

\begin{center}
\begin{tabular}{@{} l l l @{}}
\toprule
\textbf{Method} & \textbf{Half-width rate} & \textbf{Source} \\
\midrule
Fixed-sample CI & $\displaystyle\frac{\sigma}{\sqrt{n}} \cdot z_{\alpha/2}$
  & Central Limit Theorem \\
Normal mixture CS & $\displaystyle\frac{1}{n}\sqrt{(n\sigma^2 + \nu)
  \log\frac{n\sigma^2 + \nu}{\nu\alpha^2}}$ & Howard et al.\ (2021) \\
LIL-rate CS & $\displaystyle\frac{\sigma}{\sqrt{n}} \cdot
  \sqrt{2\log\log(n\sigma^2)}$ & Law of iterated logarithm \\
\bottomrule
\end{tabular}
\end{center}

The Normal mixture CS is wider than the fixed CI by a factor of
$\approx\!\sqrt{\log n}$ (the price of peeking).  It is slightly wider than the
LIL-rate CS asymptotically, but \textbf{tighter in practice} over the 2--3
orders of magnitude of sample sizes relevant to A/B testing.

\subsubsection*{Why Eppo chose this framework}

The Normal mixture boundary is the formula at the heart of Eppo's sequential
testing system (Schmit \& Miller, 2024).  To see why, consider what an A/B
testing platform needs:

\begin{enumerate}[nosep]
  \item \textbf{Anytime-valid:} Experimenters will peek --- the guarantee must
        hold at every moment.  \checkmark\ (Ville's inequality via the
        sub-$\psi$ supermartingale.)
  \item \textbf{No known variance required:} Real experiments have unknown,
        potentially changing variance.  \checkmark\ (The sub-$\psi$ framework
        permits estimated variance; Eppo plugs in $\hat{\sigma}^2(t)$ at each
        time step.)
  \item \textbf{No distributional assumption:} Metrics like revenue or session
        length are far from Normal.  \checkmark\ (Sub-Gaussian condition is
        satisfied by all bounded data and approximately by most practical
        metrics at scale.)
  \item \textbf{Tunable tightness:} Different experiments have different
        expected durations.  \checkmark\ (The parameter $\nu$ lets Eppo set
        $\nu = M \hat{\sigma}^2$, where $M$ is the expected total sample size,
        making the boundary tightest around the planned analysis time.)
  \item \textbf{Produces a confidence interval:} Experimenters need to report
        ``the lift is $+3\% \pm 1.5\%$,'' not just ``significant / not
        significant.''  \checkmark\ (The CS directly gives an interval
        $\bar{x}_t \pm u(v_t)/t$ at every time.)
\end{enumerate}

In Acts~9--12, we will see how Eppo wraps this boundary into a complete
pipeline: the confidence sequence is applied to \emph{regression-adjusted}
observations (generalised CUPED) and reported as a confidence interval around
the \textbf{relative lift} $\mu_1/\mu_0 - 1$.  The key formula in Act~10 ---
\[
  \text{CI}(t) = \hat{\tau}(t) \;\pm\; \hat{\sigma}_{\hat{\tau}}(t) \cdot
  \sqrt{\frac{n + \nu}{n} \cdot \log\frac{n + \nu}{\nu \alpha^2}}
\]
--- is just the Normal mixture boundary $u(v)$ from this act, with $v$
replaced by the estimated cumulative variance $n \hat{\sigma}_{\hat{\tau}}^2(t)$
and the boundary divided by $n$ to convert from the sum scale to the mean
scale.

\begin{keytakeaway}
\textbf{Key concepts:} confidence sequences (uniform coverage over all times),
the Normal mixture boundary $u(v) = \sqrt{(v+\nu)\log\frac{v+\nu}{\nu\alpha^2}}$,
the tuning parameter $\nu$, the sub-$\psi$ framework, nonparametric validity,
the $\sqrt{\log n}$ ``price of peeking.''

\textbf{This is the framework Eppo uses.}  The boundary formula from this act
is the engine inside Eppo's confidence interval.
\end{keytakeaway}


% ══════════════════════════════════════════════════════════════════════════════
% ACT 9
% ══════════════════════════════════════════════════════════════════════════════
\newpage
\section{Act~9 --- The Problem Eppo Solves}

\subsection{The Simulation}

\begin{simulation}
\textbf{What is being simulated:}
10,000 parallel A/B experiments, each with a small true treatment effect,
comparing classical and sequential approaches.

\textbf{Panel~1 --- Classical A/B test with peeking (failure case):}
Two groups (control / treatment) drawn from distributions with a small true
difference.  A classical 95\% CI is recomputed at each time step.

10,000 parallel experiments.  Counters show:
\begin{itemize}[nosep]
  \item ``False conclusions if you stop when CI excludes 0'': far above 5\%.
  \item ``Experiments still running after planned duration'': many finish
        early with a wrong answer.
\end{itemize}

\textbf{False decision visualisation:}
Experiments that stop early with the \emph{wrong} conclusion are highlighted
in red.  A running counter displays \textbf{``Wrong winners: $k$ / $N$''}
--- making it viscerally clear how many experiments would ship the wrong
feature.

\textbf{Panel~2 --- Sequential CI (preview):}
Same experiments, but now using a sequential confidence interval (from Act~8).
The band is wider at first, but the false conclusion rate stays $\leq 5\%$
regardless of when you peek.  Incorrectly-stopped experiments are again
highlighted, but the count stays near the $\alpha$ level.

\textbf{Panel~3 --- Time-to-detection comparison:}
Two versions of the same experiment side-by-side:
\begin{itemize}[nosep]
  \item \textbf{Without variance reduction (red):} Raw outcomes.  CI is wide.
        Takes many observations to detect a real effect.
  \item \textbf{With variance reduction (blue):} Adjusted outcomes (noise
        removed).  CI is narrow.  Same effect detected much sooner.
\end{itemize}
A counter displays: \textbf{``Average samples to detect effect: $n_{\text{raw}}$
vs.\ $n_{\text{adjusted}}$''} --- quantifying the speed-up.

\textbf{Controls:}
\begin{itemize}[nosep]
  \item Effect size slider $\delta$ (0 to 10\%).
  \item Variance level slider (low / medium / high).
  \item Number of parallel experiments (100 / 1,000 / 10,000).
  \item Fixed seed toggle (deterministic replay).
\end{itemize}

\textbf{Key takeaway:}
Classical CIs break under peeking; sequential CIs stay valid.  Variance
reduction dramatically speeds up detection.
\end{simulation}

\subsection{Intuitive Explanation}

\begin{intuition}
A real A/B testing platform faces three problems simultaneously:

\begin{enumerate}
  \item \textbf{The peeking problem} (Act~0): Product managers check results
        daily.  Classical tests break under repeated checking.

  \item \textbf{High variance:} Users are wildly different.  Some buy
        \texteuro0, others buy \texteuro500.  This noise drowns out the
        signal --- a 2\% lift buried under 200\% noise.  Experiments take
        weeks to reach significance.

  \item \textbf{Heterogeneous users:} Not all users respond the same way.
        Heavy spenders, occasional browsers, mobile vs.\ desktop --- all
        mixed together.  Covariates (pre-experiment behaviour, device type,
        country) carry information that can separate signal from noise.
\end{enumerate}

\textbf{What Eppo wants:}
\begin{itemize}[nosep]
  \item \textbf{Anytime-valid inference:} Peek as often as you like;
        the error guarantee holds.
  \item \textbf{Faster experiments:} Use pre-experiment data to reduce noise,
        shrinking the confidence interval.
  \item \textbf{Simple outputs:} Report a confidence interval for the
        \emph{relative lift} (``$+3\% \pm 1.5\%$'') that experimenters can
        act on immediately.
\end{itemize}

Acts~0--8 built each ingredient.  The next three acts show how Eppo assembles
them into one pipeline.
\end{intuition}

\begin{keytakeaway}
\textbf{Three problems, three solutions:}
\begin{center}
\renewcommand{\arraystretch}{1.3}
\begin{tabular}{@{} l l l @{}}
\toprule
\textbf{Problem} & \textbf{Solution} & \textbf{Where we learned it} \\
\midrule
Peeking inflates errors & Sequential CI (Howard et al.) & Acts~5, 8 \\
High variance / slow experiments & Regression adjustment (CUPED) & Act~11 \\
Need interpretable output & Relative lift with CI & Act~12 \\
\bottomrule
\end{tabular}
\end{center}
\end{keytakeaway}


% ══════════════════════════════════════════════════════════════════════════════
% ACT 10
% ══════════════════════════════════════════════════════════════════════════════
\newpage
\section{Act~10 --- The Eppo Pipeline, Step by Step\\
         (Schmit \& Miller, 2024)}

This is the central act of the Eppo section.  It walks through the entire
pipeline from data collection to decision, one step at a time.

\subsection{The Simulation}

\begin{simulation}
\textbf{What is being simulated:}
A single A/B experiment flowing through Eppo's pipeline, stage by stage.

\textbf{Animated pipeline diagram.}  Data flows left-to-right through
labelled stages.  Each stage lights up as it is explained.

\textbf{Stages (rendered as connected boxes):}
\[
  \fbox{Randomise}
  \;\to\;
  \fbox{Collect}
  \;\to\;
  \fbox{Adjust}
  \;\to\;
  \fbox{Estimate}
  \;\to\;
  \fbox{CI}
  \;\to\;
  \fbox{Decide}
\]

\textbf{Step-by-step mode:}
A button labelled \textbf{``Step through pipeline''} pauses at each stage.
At each stage, a distinct visual state is shown:

\begin{center}
\renewcommand{\arraystretch}{1.3}
\begin{tabular}{@{} l p{9cm} @{}}
\toprule
\textbf{Stage} & \textbf{Visual state} \\
\midrule
Randomise & Users appear as dots, randomly coloured blue (control) or
  green (treatment). \\
Collect & Each dot acquires labels: a pre-experiment value $X$ and an
  outcome $Y$.  Tooltip on hover shows both values. \\
Adjust & Dots animate: $Y \to Y^*$.  A regression line overlays each group.
  Dots shift vertically to their residual position.  Spread visibly
  decreases. \\
Estimate & Group means appear as horizontal lines.  The lift calculation
  ($\hat{\tau}_{\text{rel}}$) displays as a labelled gap between lines. \\
CI & A confidence interval band draws around the estimated lift.
  The band narrows as more data arrives. \\
Decide & The CI is colour-coded:
  \textcolor{green!60!black}{\textbf{green}} (entirely above 0),
  \textcolor{red}{\textbf{red}} (entirely below 0), or
  \textcolor{gray}{\textbf{grey}} (crosses 0).  A verdict label appears. \\
\bottomrule
\end{tabular}
\end{center}

\textbf{Toggle: ``With regression adjustment'' / ``Without regression
adjustment.''}
When toggled off, the ``Adjust'' box is greyed out, dots do not shift,
and the CI is visibly wider.

\textbf{Side-by-side comparison mode:}
An optional button labelled \textbf{``Compare with / without adjustment''}
splits the display into two panels:
\begin{itemize}[nosep]
  \item \textbf{Left panel:} no regression adjustment (raw outcomes).
  \item \textbf{Right panel:} with regression adjustment (residuals).
\end{itemize}
Both panels use the \emph{same underlying data} --- only the adjustment
differs.  The CI width difference is immediately visible, making the
benefit of variance reduction obvious at a glance.

\textbf{Stage highlights (step-by-step mode):}
When stepping through the pipeline, each stage displays a
\textbf{``Before / After''} mini-panel showing the data state entering and
leaving that stage.  Key changes are highlighted:
\begin{itemize}[nosep]
  \item \textbf{Adjust stage:} scatter shifts from raw to residual; variance
        annotation updates.
  \item \textbf{Estimate stage:} group means and lift label appear.
  \item \textbf{CI stage:} band draws with colour coding (green / red / grey).
\end{itemize}

\textbf{Controls:}
\begin{itemize}[nosep]
  \item True effect size $\delta$ (0 to 10\%).
  \item Sample size $n$ (100 to 100,000).
  \item Pre-experiment correlation $\rho$ (0 to 0.95): controls how much
        noise the adjustment removes.
  \item Fixed seed toggle (deterministic replay).
  \item ``Run 1,000 experiments'' mode (batch statistics).
\end{itemize}

\textbf{Key takeaway:}
Each pipeline stage transforms the data in a specific, visible way.
Regression adjustment visibly tightens the CI; removing it shows the cost.
\end{simulation}

\subsection{Intuitive Explanation}

\begin{intuition}
Think of Eppo's pipeline as a factory with six stations.  Each station does
one job.  Data enters raw and exits as a trustworthy confidence interval.
\end{intuition}

\subsection{The Pipeline}

\subsubsection*{Step 1: Randomise}

Each user arriving at the website is randomly assigned to \textbf{control}
(existing experience) or \textbf{treatment} (new feature).  Randomisation
ensures the groups are comparable --- any difference in outcomes is caused by
the treatment, not by pre-existing differences.

\subsubsection*{Step 2: Collect data over time}

As users interact with the product, two quantities are recorded for each user:
\begin{itemize}[nosep]
  \item $Y_i$: the \textbf{outcome} during the experiment (e.g.\ purchase
        amount, clicks, session length).
  \item $X_i$: a \textbf{pre-experiment covariate} --- the same metric from
        before the experiment started (e.g.\ purchases in the prior 2 weeks),
        plus optionally other covariates (country, device type, etc.).
\end{itemize}

Data arrives continuously.  The pipeline processes it in batches or
continuously as it flows in.

\subsubsection*{Step 3: Regression adjustment (noise removal)}

This is where Eppo \emph{generalises} the CUPED idea from Deng et al.\ (2013).

For each group \textbf{separately}:
\begin{enumerate}[nosep]
  \item Fit a regression model predicting $Y$ from the pre-experiment
        covariates $X$:
        \[
          \hat{f}_{\text{control}}(X), \qquad \hat{f}_{\text{treatment}}(X)
        \]
  \item Compute the \textbf{adjusted outcome} (residual) for each user:
        \[
          Y^*_i = Y_i - \hat{f}_g(X_i)
        \]
        where $g$ is the user's group.
\end{enumerate}

\begin{plainlanguage}
``Predict what each user would have done based on their history, then subtract
that prediction.  What's left is the part you \emph{couldn't} have predicted
--- and that's where the treatment effect hides.''
\end{plainlanguage}

\textbf{Why separate models?}  The relationship between covariates and outcome
might differ between groups (e.g.\ heavy users may respond more to the
treatment).  Fitting separately captures this.

\textbf{Why this generalises CUPED:}  Classic CUPED uses a single linear
coefficient $\theta$ applied identically to both groups.  Eppo's approach
allows nonlinear models and group-specific coefficients --- strictly more
general.

\subsubsection*{Step 4: Estimate the treatment effect}

Compute the adjusted group means:
\[
  \hat{\mu}_0(t) = \frac{1}{n_0(t)} \sum_{i \in \text{control}} Y^*_i, \qquad
  \hat{\mu}_1(t) = \frac{1}{n_1(t)} \sum_{i \in \text{treatment}} Y^*_i
\]

The estimated treatment effect:
\[
  \hat{\tau}(t) = \hat{\mu}_1(t) - \hat{\mu}_0(t)
\]

And the \textbf{relative lift} (dividing by the control mean):
\[
  \boxed{\text{Relative lift}(t) = \frac{\hat{\tau}(t)}{\hat{\mu}_0(t)}
  = \frac{\hat{\mu}_1(t)}{\hat{\mu}_0(t)} - 1}
\]

\begin{plainlanguage}
If control averages \texteuro100 and treatment averages \texteuro103, the
relative lift is $103/100 - 1 = 0.03 = +3\%$.
\end{plainlanguage}

\subsubsection*{Why relative lift?}

Eppo reports the \textbf{relative} lift $\hat{\tau}(t)/\hat{\mu}_0(t)$
rather than the absolute difference $\hat{\tau}(t)$.  Why?

\begin{center}
\renewcommand{\arraystretch}{1.3}
\begin{tabular}{@{} l c c c @{}}
\toprule
\textbf{Scenario} & \textbf{Control mean} & \textbf{Treatment mean}
  & \textbf{Lift} \\
\midrule
High-value metric   & \texteuro100 & \texteuro103 & $+3\%$ \\
Low-value metric    & \texteuro10  & \texteuro13  & $+30\%$ \\
\bottomrule
\end{tabular}
\end{center}

Both scenarios have the same absolute difference (\texteuro3), but the
business impact is very different.  Relative lift captures this:

\begin{itemize}[nosep]
  \item \textbf{Scale invariance:} A 3\% lift means the same thing whether
        the baseline is \texteuro10 or \texteuro10{,}000.
  \item \textbf{Comparability:} You can compare lift across experiments on
        different metrics (clicks, revenue, sessions) because the denominator
        normalises the scale.
  \item \textbf{Business alignment:} Product teams think in percentages ---
        ``this feature increased conversion by 2\%'' --- not in raw units.
\end{itemize}

\textbf{UI toggle:} The dashboard supports a toggle between
\textbf{``Show absolute difference''} and \textbf{``Show relative lift.''}
Both are derived from $\hat{\tau}(t)$; relative lift simply divides by the
control mean.

\subsubsection*{Step 5: Estimate the variance}

Compute the standard error of $\hat{\tau}(t)$ from the \textbf{residuals}
$Y^*_i$:
\[
  \hat{\sigma}_{\hat{\tau}}^2(t) = \frac{s_0^2(t)}{n_0(t)} + \frac{s_1^2(t)}{n_1(t)}
\]
where $s_g^2(t)$ is the sample variance of the adjusted outcomes in group~$g$.

\begin{plainlanguage}
The variance is estimated from the data --- not assumed known.  This is a key
difference from the mSPRT (Act~7), which required known variance.  Howard et
al.'s framework (Act~8) permits this.
\end{plainlanguage}

\subsubsection*{Step 6: Construct the sequential confidence interval}

Plug the estimated treatment effect and its standard error into the Normal
mixture boundary from Act~8:
\[
  \boxed{\text{CI}(t) = \hat{\tau}(t) \;\pm\;
  \hat{\sigma}_{\hat{\tau}}(t) \cdot
  \sqrt{\frac{n + \nu}{n} \cdot \log\!\frac{n + \nu}{\nu \alpha^2}}}
\]

where $n = n_0(t) + n_1(t)$ is the total sample size and $\nu = M \cdot
\hat{\sigma}^2$ is the tuning parameter (set from the expected total sample
size $M$).

\subsubsection*{Step 7: Decide}

\begin{center}
\renewcommand{\arraystretch}{1.3}
\begin{tabular}{@{} l l p{6.5cm} @{}}
\toprule
\textbf{What you see} & \textbf{Colour} & \textbf{Action} \\
\midrule
CI entirely above 0
  & \textcolor{green!60!black}{\rule{1em}{1em}} Green
  & Positive effect.  Ship the feature. \\
CI entirely below 0
  & \textcolor{red}{\rule{1em}{1em}} Red
  & Negative effect.  Revert. \\
CI crosses 0, experiment ongoing
  & \textcolor{gray}{\rule{1em}{1em}} Grey
  & Inconclusive.  Keep collecting. \\
CI crosses 0, experiment ended
  & \textcolor{gray}{\rule{1em}{1em}} Grey
  & No significant effect.  Decide pragmatically. \\
\bottomrule
\end{tabular}
\end{center}

Because the CI is a confidence sequence (Act~8), \textbf{this decision is
valid at whatever time it is made.}  No need to pre-specify an analysis time.

\begin{keytakeaway}
\textbf{The Eppo pipeline in one sentence:}
Randomise $\to$ collect $Y$ and covariates $X$ $\to$ regression-adjust
(remove predictable noise) $\to$ estimate lift and variance $\to$ wrap in
a sequential CI $\to$ decide when CI excludes zero.

Every step matters.  The regression removes noise (Act~11 dives deeper).
The sequential CI guarantees validity (Act~8).  The relative lift provides
the business answer.
\end{keytakeaway}


% ══════════════════════════════════════════════════════════════════════════════
% ACT 11
% ══════════════════════════════════════════════════════════════════════════════
\newpage
\section{Act~11 --- Variance Reduction via Regression Adjustment}

\subsection{The Simulation}

\begin{simulation}
\textbf{What is being simulated:}
The effect of regression adjustment on variance and confidence interval width,
comparing raw outcomes with adjusted outcomes.

\textbf{Panel~1 --- Why raw averages are noisy:}
Scatter plot of 1,000 users: $x$-axis = pre-experiment purchases,
$y$-axis = experiment purchases.  Strong positive correlation visible.
Two horizontal lines show the raw group means (control in blue, treatment
in green).  The gap between them \emph{is} the treatment effect --- but the
scatter around each line is enormous.

\textbf{Panel~2 --- Regression adjustment:}
Same scatter plot, but now a regression line is fitted per group.  Each dot
is replaced by its \textbf{residual} (vertical distance from the regression
line).  The residuals have much less spread.  The adjusted group means are
shown --- same gap, much less scatter.

\textbf{Panel~3 --- Effect on the confidence interval:}
Two confidence bands evolving over time:
\begin{itemize}[nosep]
  \item \textbf{Red (raw):} Wide band.  Slowly narrows.
  \item \textbf{Blue (adjusted):} Narrow band.  Reaches significance much
        sooner.
\end{itemize}
A counter shows: ``Observations saved: $X\%$'' --- how many fewer
observations the adjusted version needs to reach the same width.

\textbf{Controls:}
\begin{itemize}[nosep]
  \item Correlation slider $\rho$ (0 to 0.95): controls how predictable
        users are from their pre-experiment behaviour.
  \item True effect size $\delta$ (0 to 10\%).
  \item Toggle: with/without adjustment.
  \item Fixed seed toggle (deterministic replay).
  \item ``Run 1,000 experiments'' mode (batch statistics).
\end{itemize}

\textbf{Key takeaway:}
The higher the correlation $\rho$ between pre-experiment and experiment data,
the more variance is removed.  At $\rho = 0.7$, the experiment runs
$\approx 2\times$ faster.
\end{simulation}

\subsection{Intuitive Explanation}

\begin{intuition}
Imagine weighing yourself before and after a meal.  If every person in the
study has a different baseline weight (50\,kg to 120\,kg), the raw
``post-meal weight'' is extremely variable.  But the \emph{change} (before
$\to$ after) is much more consistent --- a meal adds roughly the same weight
regardless of the person.

Regression adjustment does the same thing.  Each user's outcome $Y$ is
partly explained by their pre-experiment behaviour $X$ (their ``baseline
weight'').  The regression model predicts that baseline.  Subtracting the
prediction leaves the \emph{change} --- the part due to the treatment, not
the person.

\textbf{The result:} less noise, same signal, tighter confidence intervals,
faster decisions.
\end{intuition}

\subsection{Mathematical Formulation}

\subsubsection*{From CUPED to generalised regression adjustment}

Recall CUPED (Deng et al., 2013) from the notation in Act~10, Step~3:
\[
  Y^* = Y - \theta(X - \E[X])
\]
with optimal $\theta = \Cov(Y,X)/\Var(X)$ and variance reduction
$\Var(Y^*) = \Var(Y)(1 - \rho^2)$, where $\rho = \Cor(Y,X)$ is the
\textbf{correlation coefficient} between the outcome and the covariate
(a number between $-1$ and $+1$ measuring how strongly they move together).

Eppo generalises this in two ways:
\begin{enumerate}[nosep]
  \item \textbf{Nonlinear models:}  Instead of a single linear coefficient
        $\theta$, fit a full regression model $\hat{f}(X)$ that can capture
        nonlinear relationships (e.g.\ random forests, gradient boosting).
  \item \textbf{Separate models per group:}  Fit $\hat{f}_{\text{control}}(X)$
        and $\hat{f}_{\text{treatment}}(X)$ independently.  This allows the
        treatment to change how covariates relate to outcomes.
\end{enumerate}

The adjusted outcome for user~$i$ in group~$g$:
\[
  \boxed{Y^*_i = Y_i - \hat{f}_g(X_i)}
\]

\begin{plainlanguage}
``Take what user~$i$ actually did.  Subtract what the model predicted they
would do based on their history.  What's left is the residual --- the
\emph{unpredictable} part.''
\end{plainlanguage}

\subsubsection*{Why this reduces variance}

The variance of the residual $Y^*$ is:
\[
  \Var(Y^*) = \Var(Y - \hat{f}(X)) = \Var(Y)(1 - R^2)
\]
where $R^2$ is the coefficient of determination of the regression model ---
the fraction of variance in $Y$ that is explained by $X$.

\begin{center}
\begin{tabular}{@{} c c r l @{}}
\toprule
$R^2$ & $1-R^2$
       & \textbf{Variance remaining} & \textbf{Equivalent traffic multiplier} \\
\midrule
0.10  & 0.90 & 90\%  & $1.1\times$ \\
0.25  & 0.75 & 75\%  & $1.3\times$ \\
0.50  & 0.50 & 50\%  & $2\times$ \\
0.60  & 0.40 & 40\%  & $2.5\times$ \\
0.75  & 0.25 & 25\%  & $4\times$ \\
\bottomrule
\end{tabular}
\end{center}

\begin{plainlanguage}
With $R^2 = 0.50$, the regression explains half the variance.  The
remaining noise is halved --- equivalent to running the experiment with twice
as many users.
\end{plainlanguage}

\subsubsection*{What covariate to use}

Following Deng et al.\ (2013):
\begin{enumerate}[nosep]
  \item \textbf{The same metric from the pre-experiment period} is
        consistently the most powerful covariate.
  \item A pre-period of \textbf{$\sim\!2$ weeks} works well for experiments
        lasting 1--3 weeks (longer pre-periods have diminishing returns).
  \item Additional covariates (device type, country, day-of-week of first
        visit) provide marginal gains.
\end{enumerate}

\textbf{Practical results} (Deng et al., 2013, at Microsoft Bing):
\begin{itemize}[nosep]
  \item Queries-per-user, clicks-per-user: \textbf{45--52\%} variance
        reduction ($\Rightarrow$ roughly $2\times$ traffic multiplier).
  \item Revenue-per-user: $< 5\%$ reduction (low temporal correlation).
\end{itemize}

\subsubsection*{The critical safety rule}

\begin{keytakeaway}
The covariate $X$ must satisfy $\E[X^{(\text{treatment})}] =
\E[X^{(\text{control})}]$.  This is automatically guaranteed when $X$ comes
from the \emph{pre-experiment} period (randomisation has not yet occurred).

\textbf{NEVER use post-treatment covariates.}  Deng et al.\ demonstrated
that using an in-experiment covariate (highly correlated with the outcome)
produced \textbf{directionally opposite conclusions}, because the treatment
itself changed $X$.
\end{keytakeaway}

\subsubsection*{Why adjustment does not break the sequential guarantee}

The confidence sequence (Act~8) is applied to the adjusted outcomes $Y^*$,
not the raw $Y$.  The guarantee holds because:
\begin{enumerate}[nosep]
  \item Covariates come from the pre-experiment period (independent of
        treatment).
  \item Regression adjustment changes only the variance, not the mean
        treatment effect.
  \item The sub-Gaussian condition (Act~8) is satisfied by the adjusted
        residuals.
  \item Ville's inequality applies to the resulting supermartingale.
\end{enumerate}

\begin{plainlanguage}
The adjustment is like using a better measuring instrument --- more precise,
but measuring the same thing.  The sequential guarantee depends on the
structure of the martingale, which is preserved.
\end{plainlanguage}

\begin{keytakeaway}
\textbf{Key concepts:} regression adjustment as generalised CUPED, separate
models per group, variance reduction factor $1 - R^2$, the safety rule
(pre-experiment covariates only), adjustment preserves the sequential
guarantee.
\end{keytakeaway}


% ══════════════════════════════════════════════════════════════════════════════
% ACT 12
% ══════════════════════════════════════════════════════════════════════════════
\newpage
\section{Act~12 --- Sequential Confidence Intervals in Eppo}

\subsection{The Simulation}

\begin{simulation}
\textbf{What is being simulated:}
The behaviour of Eppo's sequential confidence interval versus a classical CI,
with controls for the tuning parameter $\nu$ and regression adjustment.

\textbf{Panel~1 --- Classical vs.\ sequential CI over time:}
A single experiment with a true effect of $+3\%$.  Two bands evolve over time:
\begin{itemize}[nosep]
  \item \textbf{Dashed grey:} Classical 95\% CI (recomputed at each step).
        Narrow but invalid under peeking.
  \item \textbf{Solid blue:} Eppo's sequential CI.  Wider at first, then
        approaches the classical CI as $n$ grows.
\end{itemize}

A counter shows the ``sequential multiplier'' --- the factor that replaces
$z_{\alpha/2} \approx 1.96$ in the classical CI.  It starts high (e.g.\ 3.5)
and slowly decreases.

\textbf{Panel~2 --- $\nu$ tuning:}
A slider for $\nu$ (or equivalently $M$, the expected total sample size).
\begin{itemize}[nosep]
  \item Small $\nu$ (small $M$): CI is tight early, wide late.
  \item Large $\nu$ (large $M$): CI is wide early, tight late.
  \item $\nu \approx M\hat{\sigma}^2$: sweet spot at the planned analysis
        time.
\end{itemize}

An animation shows the boundary curve $u(v)$ reshaping as $\nu$ changes.

\textbf{Panel~3 --- 10,000 experiments:}
Run 10,000 parallel experiments.  Two counters:
\begin{itemize}[nosep]
  \item \textbf{Classical CI with peeking:} ``Fraction that ever falsely
        excluded 0'': $\gg 5\%$.
  \item \textbf{Sequential CI with peeking:} Same counter: $\leq 5\%$.
\end{itemize}

\textbf{Controls:}
\begin{itemize}[nosep]
  \item True effect size $\delta$ (0 for null, or a positive value).
  \item Expected sample size $M$ (slider that sets $\nu$).
  \item Significance level $\alpha$.
  \item Toggle: with/without regression adjustment.
  \item Fixed seed toggle (deterministic replay).
  \item ``Run 1,000 experiments'' mode (batch statistics).
\end{itemize}

\textbf{Key takeaway:}
The sequential CI is wider than the classical CI (the ``price of peeking'')
but guarantees validity under continuous monitoring.  The tuning parameter
$\nu$ controls where the CI is tightest.
\end{simulation}

\subsection{Intuitive Explanation}

\begin{intuition}
This act connects the pipeline from Act~10 to the final statistical output
that an experimenter sees on their dashboard.

A standard 95\% CI says: ``Right now, the lift is $+3\% \pm 1.5\%$.''  But
that guarantee is only valid at this exact moment.  If you checked yesterday,
and you check again tomorrow, the guarantee has already been violated.

Eppo's sequential CI says: ``The lift is $+3\% \pm 2.1\%$, and this has been
valid \emph{at every moment since the experiment started}.''

The price: the interval is wider --- the ``$\pm 2.1\%$'' instead of
``$\pm 1.5\%$''.  This is the \textbf{cost of peeking}.  But it is a small
price for the freedom to check whenever you want and stop as soon as you are
convinced.
\end{intuition}

\subsubsection*{Why the interval shrinks slower}

A standard CI shrinks as $1/\sqrt{n}$: double your sample, the interval
narrows by $\sqrt{2} \approx 1.41$.  Eppo's sequential CI shrinks slower ---
roughly as $\sqrt{\log n / n}$ --- because it must remain valid across
\emph{all} possible stopping times, not just the one you chose.

\begin{plainlanguage}
Think of a standard CI as a promise valid at one specific moment: ``At
$n = 10{,}000$, I'm 95\% confident.''  A sequential CI makes a much
stronger promise: ``At every $n$ from 1 to $\infty$, I'm 95\% confident.''

The extra $\sqrt{\log n}$ factor is the price of that stronger promise.
As data accumulates, the penalty grows so slowly that it barely matters ---
$\sqrt{\log(10{,}000)} \approx 3.0$, while $\sqrt{\log(100{,}000)} \approx
3.4$ --- but it never disappears entirely.
\end{plainlanguage}

\subsection{Mathematical Formulation}

\subsubsection*{The formula, revisited}

From the pipeline (Act~10, Step~6):
\[
  \boxed{\text{CI}(t) = \hat{\tau}(t) \;\pm\;
  \hat{\sigma}_{\hat{\tau}}(t) \cdot
  \underbrace{\sqrt{\frac{n + \nu}{n} \cdot
  \log\!\frac{n + \nu}{\nu \alpha^2}}}_{\text{sequential multiplier}}}
\]

Compare with a standard CI:
\[
  \text{CI}_{\text{classical}}(t) = \hat{\tau}(t) \;\pm\;
  \hat{\sigma}_{\hat{\tau}}(t) \cdot
  \underbrace{z_{\alpha/2}}_{\approx\,1.96}
\]

The only difference is the \textbf{sequential multiplier} replacing the
constant $1.96$.

\subsubsection*{How the sequential multiplier behaves}

Define:
\[
  m(n) = \sqrt{\frac{n + \nu}{n} \cdot \log\frac{n + \nu}{\nu \alpha^2}}
\]

\begin{center}
\renewcommand{\arraystretch}{1.3}
\begin{tabular}{@{} r c c l @{}}
\toprule
$n$ & $m(n)$ (example) & $z_{\alpha/2}$ & \textbf{Interpretation} \\
\midrule
100 & $\approx 3.8$ & 1.96 & Early: sequential CI is $\sim\!2\times$ wider \\
1{,}000 & $\approx 2.8$ & 1.96 & Mid-experiment: gap narrowing \\
10{,}000 & $\approx 2.3$ & 1.96 & Late: only $\sim\!17\%$ wider \\
100{,}000 & $\approx 2.1$ & 1.96 & Approaching classical, never quite reaching it \\
\bottomrule
\end{tabular}
\end{center}

(Exact values depend on $\nu$, $\alpha$; the above is illustrative for
$\nu = 10{,}000$, $\alpha = 0.05$.)

\begin{plainlanguage}
The sequential multiplier starts high and shrinks toward $z_{\alpha/2}$ as
more data arrives.  The gap --- the ``price of peeking'' --- decreases but
never fully closes.  For most practical experiments (thousands to hundreds of
thousands of observations), the price is modest: 10--40\% wider than a
classical CI.
\end{plainlanguage}

\subsubsection*{The role of $\nu$ and $M$}

Recall from Act~8: $\nu$ controls \emph{which sample size} the boundary is
tightest for.  Schmit \& Miller set:
\[
  \nu = M \cdot \hat{\sigma}^2
\]
where $M$ is the \textbf{expected total sample size} (planned experiment
duration $\times$ daily traffic).

\begin{center}
\begin{tabular}{@{} l l @{}}
\toprule
\textbf{If\ldots} & \textbf{Then\ldots} \\
\midrule
$M$ is too small (underestimate traffic) & CI is tight early but wide late \\
$M$ is too large (overestimate traffic) & CI is wide early but tight late \\
$M$ matches actual traffic & CI is tightest around the planned end date \\
\bottomrule
\end{tabular}
\end{center}

\begin{plainlanguage}
Setting $M$ is like aiming a spotlight: you choose where the beam is brightest.
Aim at the expected experiment end, and you get the tightest CI right when you
need to make the decision.  If the experiment runs longer or shorter than
planned, the CI is still valid --- just not optimally tight.
\end{plainlanguage}

\subsubsection*{Reporting: from treatment effect to relative lift}

Eppo reports the relative lift and its CI:
\[
  \text{Relative lift}(t) = \frac{\hat{\tau}(t)}{\bar{Y}^*_{\text{control}}(t)},
  \qquad
  \text{CI}_{\text{lift}}(t) = \frac{\text{CI}(t)}{\bar{Y}^*_{\text{control}}(t)}
\]

\begin{plainlanguage}
Dividing by the control mean converts the absolute difference into a
percentage.  ``The treatment increased purchases by \texteuro3 per user''
becomes ``the treatment increased purchases by $+3\%$.''
\end{plainlanguage}

\subsubsection*{Decision rule}

\begin{center}
\renewcommand{\arraystretch}{1.3}
\begin{tabular}{@{} l l p{6.5cm} @{}}
\toprule
\textbf{What you see} & \textbf{Colour} & \textbf{What it means} \\
\midrule
CI entirely above 0
  & \textcolor{green!60!black}{\rule{1em}{1em}} Green
  & Positive effect.  Ship it. \\
CI entirely below 0
  & \textcolor{red}{\rule{1em}{1em}} Red
  & Negative effect.  Revert. \\
CI crosses 0, experiment ongoing
  & \textcolor{gray}{\rule{1em}{1em}} Grey
  & Inconclusive.  Keep collecting. \\
CI crosses 0, experiment ended
  & \textcolor{gray}{\rule{1em}{1em}} Grey
  & No significant effect.  Decide pragmatically. \\
\bottomrule
\end{tabular}
\end{center}

\textbf{At any of these moments, the decision is valid.}  This is the payoff
of the entire mathematical journey from Act~0: Ville's inequality, applied
through eight acts of increasingly practical machinery, guarantees that the
false positive rate stays $\leq \alpha$ no matter when you look.

\begin{keytakeaway}
\textbf{What Eppo actually uses (summary):}
\begin{itemize}[nosep]
  \item \textbf{Framework:} Howard et al.\ (2021) confidence sequences
        (not the mSPRT of Johari et al.\ 2017).
  \item \textbf{Boundary:} The Normal mixture boundary
        $u(v) = \sqrt{(v+\nu)\log\frac{v+\nu}{\nu\alpha^2}}$.
  \item \textbf{Tuning:} $\nu = M \hat{\sigma}^2$, where $M$ is the expected
        total sample size.
  \item \textbf{Variance:} Estimated from data (not assumed known).
  \item \textbf{Noise reduction:} Generalised CUPED via per-group
        regression adjustments.
  \item \textbf{Estimand:} Relative lift $\mu_1/\mu_0 - 1$.
  \item \textbf{Decision:} Stop when CI no longer crosses zero, valid at any
        time.
\end{itemize}
\end{keytakeaway}


% ══════════════════════════════════════════════════════════════════════════════
% ACT 13
% ══════════════════════════════════════════════════════════════════════════════
\newpage
\section{Act~13 --- What If You Don't Have Eppo?\\
         (The Poor Man's Sequential Test)}

Acts~9--12 described Eppo's sophisticated pipeline.  But many teams run
experiments without a dedicated platform.  This act presents three
progressively better approaches to peek-safe testing that anyone can implement
with a spreadsheet or a few lines of code --- and compares them head-to-head
with Eppo's confidence intervals.

\subsection{The Simulation}

\begin{simulation}
\textbf{What is being simulated:}
A single A/B experiment monitored at $K$ pre-planned analysis times, showing
four methods side-by-side.

\textbf{Panel~1 --- Four confidence bands:}
The same experiment, with four CI bands evolving over time:
\begin{itemize}[nosep]
  \item \textcolor{red}{\textbf{Red (Bonferroni):}} Widest.  Simple but
        conservative.
  \item \textcolor{orange}{\textbf{Orange (Pocock):}} Constant threshold,
        tighter than Bonferroni.
  \item \textcolor{blue}{\textbf{Blue (O'Brien--Fleming):}} Starts very wide,
        narrows to near-classical at the final analysis.
  \item \textcolor{green!60!black}{\textbf{Green (Eppo / Howard):}}
        Continuous monitoring.  No need to pre-specify $K$.
\end{itemize}

Vertical dashed lines mark the $K$ planned analysis times.  Between analyses,
the group-sequential methods (Bonferroni, Pocock, OBF) have no valid CI ---
only Eppo's band is defined everywhere.

\textbf{Panel~2 --- 10{,}000 experiments:}
Run 10{,}000 experiments under the null ($\delta = 0$).  Counters show the
false positive rate for each method:
\begin{itemize}[nosep]
  \item Classical CI with $K$ peeks: $\gg 5\%$.
  \item Bonferroni: $\leq 5\%$ (often well below).
  \item Pocock: $\leq 5\%$ (closer to 5\%).
  \item O'Brien--Fleming: $\leq 5\%$.
  \item Eppo: $\leq 5\%$.
\end{itemize}

\textbf{Panel~3 --- Efficiency comparison:}
Under a true effect ($\delta > 0$), show average stopping time (number of
observations to reach significance) for each method.

\textbf{Controls:}
\begin{itemize}[nosep]
  \item Number of planned peeks $K$ (2 to 20).
  \item True effect size $\delta$ (0 to 10\%).
  \item Significance level $\alpha$.
  \item Fixed seed toggle (deterministic replay).
  \item ``Run 1{,}000 experiments'' mode (batch statistics).
\end{itemize}

\textbf{Key takeaway:}
All four methods control the false positive rate.  The price of simplicity
is wider intervals (slower detection).  Eppo's approach is the most
efficient, but even the simplest method (Bonferroni) is far better than
naive peeking.
\end{simulation}


\subsection{Intuitive Explanation}

\begin{intuition}
You want to check your experiment every Monday.  If you use a standard
95\% CI each time, you'll get false positives far too often (Act~0).

The simplest fix: \textbf{be stricter at each check.}  If you plan to peek
4~times, demand stronger evidence each time.  The three methods below
differ in \emph{how} they distribute that strictness across the peeks.

Think of it as a budget.  You have a total error budget of $\alpha = 5\%$.
\begin{itemize}[nosep]
  \item \textbf{Bonferroni} distributes the budget equally:
        $\alpha/K$ per peek.
  \item \textbf{Pocock} distributes it more cleverly (using the correlation
        between test statistics at different peeks), but still equally across
        peeks.
  \item \textbf{O'Brien--Fleming} front-loads the budget: almost no spending
        early, most of it saved for the final analysis.
\end{itemize}

Eppo's approach (Acts~8--12) doesn't need a fixed budget at all --- it works
for \emph{any} number of peeks, continuously.
\end{intuition}


\subsection{Method 1: Bonferroni Correction (Simplest)}

\subsubsection*{The idea}

Before the experiment starts, commit to checking exactly $K$ times
(e.g.\ every Monday for 4 weeks: $K = 4$).  At each check, use a
significance level of $\alpha/K$ instead of $\alpha$.

\subsubsection*{The recipe}

\begin{enumerate}[nosep]
  \item Choose $K$ (number of planned peeks) and $\alpha$ (e.g.\ 0.05).
  \item At each peek $k = 1, \ldots, K$, compute the standard CI using
        the adjusted significance level:
        \[
          \text{CI}_k = \hat{\tau}(t_k) \;\pm\;
          \hat{\sigma}_{\hat{\tau}}(t_k) \cdot z_{\alpha/(2K)}
        \]
  \item If $0 \notin \text{CI}_k$: stop and declare significance.
  \item If $k = K$ and $0 \in \text{CI}_K$: no significant effect.
\end{enumerate}

\begin{center}
\renewcommand{\arraystretch}{1.3}
\begin{tabular}{@{} c c c c @{}}
\toprule
$K$ (peeks) & $\alpha/K$ & $z_{\alpha/(2K)}$ & CI multiplier vs.\ 1.96 \\
\midrule
1  & 0.050  & 1.96 & $1.00\times$ (no correction) \\
2  & 0.025  & 2.24 & $1.14\times$ \\
4  & 0.0125 & 2.50 & $1.28\times$ \\
8  & 0.00625 & 2.73 & $1.39\times$ \\
12 & 0.00417 & 2.87 & $1.46\times$ \\
20 & 0.0025  & 3.02 & $1.54\times$ \\
52 & 0.00096 & 3.29 & $1.68\times$ (weekly for a year) \\
\bottomrule
\end{tabular}
\end{center}

\begin{plainlanguage}
With 4 planned peeks, replace $z = 1.96$ with $z = 2.50$.  Your CI is
28\% wider at each peek, but you can safely check 4 times without
inflating errors.
\end{plainlanguage}

\subsubsection*{Pros and cons}

\begin{itemize}[nosep]
  \item[\textcolor{green!60!black}{+}] \textbf{Dead simple.}  One line of
        code: change \texttt{1.96} to \texttt{z\_alpha/(2K)}.
  \item[\textcolor{green!60!black}{+}] \textbf{Always valid.}  Works for
        any test statistic, any distribution.
  \item[\textcolor{red}{--}] \textbf{Conservative.}  The CI is wider than
        necessary because Bonferroni ignores the correlation between test
        statistics at successive peeks.
  \item[\textcolor{red}{--}] \textbf{Scales poorly.}  As $K$ grows, the
        correction becomes increasingly harsh.
\end{itemize}


\subsection{Method 2: Pocock Boundaries (1977)}

\subsubsection*{The idea}

Like Bonferroni, pre-specify $K$ equally-spaced analysis times.  But instead
of Bonferroni's crude $\alpha/K$ split, use a constant critical value $c_P$
that accounts for the correlation between test statistics at different peeks.

Because the test statistics $Z_1, Z_2, \ldots, Z_K$ at successive analyses
share overlapping data, they are positively correlated.  Pocock's method
exploits this to derive a tighter threshold.

\subsubsection*{The recipe}

\begin{enumerate}[nosep]
  \item Choose $K$ and $\alpha$.
  \item Look up (or compute) the Pocock critical value $c_P(K, \alpha)$.
  \item At each peek $k$, compute the $z$-statistic:
        $Z_k = \hat{\tau}(t_k) / \hat{\sigma}_{\hat{\tau}}(t_k)$.
  \item If $|Z_k| > c_P$: stop and declare significance.
\end{enumerate}

\begin{center}
\renewcommand{\arraystretch}{1.3}
\begin{tabular}{@{} c c c c @{}}
\toprule
$K$ & $c_P$ (Pocock) & $z_{\alpha/(2K)}$ (Bonf.) & Improvement \\
\midrule
2  & 2.18 & 2.24 & 3\% tighter \\
4  & 2.36 & 2.50 & 6\% tighter \\
8  & 2.51 & 2.73 & 8\% tighter \\
12 & 2.60 & 2.87 & 9\% tighter \\
20 & 2.69 & 3.02 & 11\% tighter \\
\bottomrule
\end{tabular}
\end{center}

\begin{plainlanguage}
Pocock gives a constant threshold that is a few percent lower than
Bonferroni's.  The saving is small when $K$ is small (3\% for 2 peeks)
but grows as the number of peeks increases (11\% for 20 peeks).

The key advantage: the threshold is the same at every peek, so it is
almost as simple to use as Bonferroni.
\end{plainlanguage}

\subsubsection*{Pros and cons}

\begin{itemize}[nosep]
  \item[\textcolor{green!60!black}{+}] \textbf{Tighter than Bonferroni.}
  \item[\textcolor{green!60!black}{+}] \textbf{Constant threshold} --- easy
        to remember and apply.
  \item[\textcolor{red}{--}] \textbf{Requires equally-spaced analyses.}
  \item[\textcolor{red}{--}] \textbf{Spends too much $\alpha$ early.}
        Because the same threshold is used early and late, it is harder to
        reject at the planned final analysis than a classical test.
\end{itemize}


\subsection{Method 3: O'Brien--Fleming Boundaries (1979)}

\subsubsection*{The idea}

Instead of a constant threshold, use a threshold that \emph{decreases} over
time.  At the first peek, demand very strong evidence; at the last peek,
demand nearly the same evidence as a classical test.

\subsubsection*{The recipe}

\begin{enumerate}[nosep]
  \item Choose $K$ and $\alpha$.
  \item At peek $k$ (of $K$ total), use the critical value:
        \[
          c_k^{\text{OBF}} \;\approx\; z_{\alpha/2} \cdot \sqrt{K/k}
        \]
        (This is the asymptotic approximation.  Exact values are computed
        numerically.)
  \item If $|Z_k| > c_k^{\text{OBF}}$: stop and declare significance.
\end{enumerate}

\begin{center}
\renewcommand{\arraystretch}{1.3}
\begin{tabular}{@{} c c c c c @{}}
\toprule
Peek $k$ (of $K{=}4$) & $c_k^{\text{OBF}}$ & $c_P$ (Pocock) &
  $z_{\alpha/(2K)}$ (Bonf.) & $z_{\alpha/2}$ (classical) \\
\midrule
1 (25\% of data) & 4.05 & 2.36 & 2.50 & 1.96 \\
2 (50\% of data) & 2.86 & 2.36 & 2.50 & 1.96 \\
3 (75\% of data) & 2.34 & 2.36 & 2.50 & 1.96 \\
4 (100\% of data) & 2.02 & 2.36 & 2.50 & 1.96 \\
\bottomrule
\end{tabular}
\end{center}

\begin{plainlanguage}
O'Brien--Fleming is extremely conservative early (threshold 4.05 at the
first peek --- you'd almost never stop that early) but barely penalises
the final analysis (2.02 vs.\ 1.96 --- only 3\% wider).

This is the most efficient approach: if you run the experiment to completion,
you pay almost nothing for the option to peek.  The cost is that early
stopping is unlikely unless the effect is very large.
\end{plainlanguage}

\subsubsection*{Pros and cons}

\begin{itemize}[nosep]
  \item[\textcolor{green!60!black}{+}] \textbf{Most efficient} of the three
        simple methods.  Final-analysis CI is barely wider than classical.
  \item[\textcolor{green!60!black}{+}] \textbf{Catches large effects early}
        (if the effect is huge, even the high early threshold is exceeded).
  \item[\textcolor{red}{--}] \textbf{Varying threshold} --- slightly harder
        to implement and communicate.
  \item[\textcolor{red}{--}] \textbf{Requires equally-spaced analyses.}
\end{itemize}


\subsection{Head-to-Head Comparison}

\begin{center}
\renewcommand{\arraystretch}{1.3}
\begin{tabular}{@{} l c c c c @{}}
\toprule
& \textbf{Bonferroni} & \textbf{Pocock} & \textbf{OBF} & \textbf{Eppo} \\
\midrule
Pre-specify $K$?
  & Yes & Yes & Yes & No \\
Threshold type
  & Constant & Constant & Decreasing & Continuous \\
CI at peek 1 ($K{=}4$)
  & $2.50\times\hat{\sigma}$ & $2.36\times\hat{\sigma}$
  & $4.05\times\hat{\sigma}$ & $\sim\!3.8\times\hat{\sigma}$ \\
CI at final analysis ($K{=}4$)
  & $2.50\times\hat{\sigma}$ & $2.36\times\hat{\sigma}$
  & $2.02\times\hat{\sigma}$ & $\sim\!2.3\times\hat{\sigma}$ \\
Valid between peeks?
  & No & No & No & Yes \\
Variance reduction?
  & Manual & Manual & Manual & Built-in \\
Implementation effort
  & 1 line & Lookup table & Formula & Platform \\
\bottomrule
\end{tabular}
\end{center}

\begin{keytakeaway}
\textbf{Which method should you use?}
\begin{itemize}[nosep]
  \item \textbf{No platform, minimal effort:} Bonferroni.  Commit to $K$
        peeks.  Replace 1.96 with $z_{\alpha/(2K)}$.  Done.
  \item \textbf{No platform, moderate effort:} O'Brien--Fleming.  Nearly as
        efficient as Eppo at the final analysis, but requires pre-specifying
        $K$ and cannot peek between scheduled analyses.
  \item \textbf{Need continuous monitoring or variance reduction:} Use a
        platform like Eppo.  The machinery of Acts~8--12 is what you're
        paying for --- continuous validity, regression adjustment, and
        no need to pre-commit to $K$.
\end{itemize}

\textbf{The bottom line:} even the crudest correction (Bonferroni) is
\emph{vastly} better than naive peeking.  If you are checking results
without any correction, you should start with Bonferroni \emph{today}.
\end{keytakeaway}

\begin{historynote}
\textbf{Stuart Pocock} (1977) and \textbf{Peter O'Brien \& Thomas Fleming}
(1979) developed these group sequential methods for clinical trials, where
stopping a trial early can save lives.  Their work sits between Wald's SPRT
(1945) and Howard et al.'s continuous framework (2021) in the historical
timeline.  The key insight of all three: if you want to peek, you must pay
--- but the price is negotiable.
\end{historynote}

\subsection{Between Acts --- Transition Cards}

Each transition features a brief interstitial card with a bridge sentence and
a ``Key insight so far'' box.

\begin{center}
\renewcommand{\arraystretch}{1.3}
\begin{tabular}{@{} l p{9cm} @{}}
\toprule
\textbf{Transition} & \textbf{Bridge sentence} \\
\midrule
Act~0 $\to$ 1 & ``To build a peek-safe test, we first need the right
  building block.  It starts with a coin flip.'' \\
Act~1 $\to$ 2 & ``Now let's put money on each step --- and learn why no
  strategy can beat a fair game.'' \\
Act~2 $\to$ 3 & ``Fair games can't be beaten.  Now let's build a detective
  --- a tool for telling two coins apart.'' \\
Act~3 $\to$ 4 & ``We can measure evidence with a likelihood ratio.  But
  here's the surprise: that evidence score is itself a martingale.'' \\
Act~4 $\to$ 5 & ``The likelihood ratio is a martingale.  Now let's see why
  that single fact gives us the peeking guarantee.'' \\
Act~5 $\to$ 6 & ``We have the theory: peeking is safe.  Now let's turn it
  into a practical decision procedure.'' \\
Act~6 $\to$ 7 & ``SPRT is optimal but inflexible.  What if we don't know
  the effect size?  Robbins had a radical idea.'' \\
Act~7 $\to$ 8 & ``The mixture approach works, but can we make it
  nonparametric?  Howard et al.\ found the answer.'' \\
Act~8 $\to$ 9 & ``The theory is complete.  Now let's see what
  problems a real A/B testing platform must solve.'' \\
Act~9 $\to$ 10 & ``Three problems, three solutions.  Now let's walk through
  Eppo's pipeline, step by step.'' \\
Act~10 $\to$ 11 & ``The pipeline uses regression adjustment to remove noise.
  Let's see exactly how and why that works.'' \\
Act~11 $\to$ 12 & ``The noise is reduced.  Now let's wrap everything in a
  sequential confidence interval and make the final decision.'' \\
Act~12 $\to$ 13 & ``Eppo solves everything --- but what if you don't have
  Eppo?  Three simple alternatives, from crude to clever.'' \\
\bottomrule
\end{tabular}
\end{center}

\subsection{Persistent Controls}

A slider panel available throughout all acts:
\begin{itemize}[nosep]
  \item Probability $p$ (or effect size $\delta$ from Act~3 onward)
  \item Significance level $\alpha$
  \item Number of steps / observations $n$
  \item Simulation speed
  \item Number of simultaneous runs
\end{itemize}

\subsection{Persistent Sidebar}

A collapsible sidebar shows the act list and current position, allowing
quick navigation to any act.

\subsection{Visual Design Principles}

\begin{itemize}[nosep]
  \item The same visual object (path / process) carries through all fourteen
        acts, relabelled at each stage.
  \item Threshold lines appear gradually: none in Acts~0--1; $\pm\sqrt{n}$
        envelope in Act~1; both SPRT thresholds in Act~6; only the upper
        threshold in Act~7; confidence bands from Act~8 onward.
  \item Formula panels slide in/out; not shown by default on first reading.
  \item Every formula has a plain-language translation directly below it.
  \item Key terms on first use are bolded and defined in-line.
\end{itemize}

\subsection{Colour Scheme}

A consistent colour palette is used across all acts:

\begin{center}
\renewcommand{\arraystretch}{1.3}
\begin{tabular}{@{} l l l @{}}
\toprule
\textbf{Element} & \textbf{Colour} & \textbf{Usage} \\
\midrule
Control group & \textcolor{blue}{\rule{1em}{1em}} Blue
  & Dots, means, CI bands for the control arm \\
Treatment group & \textcolor{green!60!black}{\rule{1em}{1em}} Green
  & Dots, means, CI bands for the treatment arm \\
CI --- positive effect & \textcolor{green!60!black}{\rule{1em}{1em}} Green
  & CI entirely above 0: ship the feature \\
CI --- negative effect & \textcolor{red}{\rule{1em}{1em}} Red
  & CI entirely below 0: revert \\
CI --- inconclusive & \textcolor{gray}{\rule{1em}{1em}} Grey
  & CI crosses 0: keep collecting \\
Classical CI (invalid) & \textcolor{gray}{\rule{1em}{1em}} Dashed grey
  & Shown for comparison; not valid under peeking \\
Raw outcomes (unadjusted) & \textcolor{red!70!black}{\rule{1em}{1em}} Red
  & Acts~11--12: unadjusted CI / scatter \\
Adjusted outcomes & \textcolor{blue}{\rule{1em}{1em}} Blue
  & Acts~11--12: regression-adjusted CI / scatter \\
\bottomrule
\end{tabular}
\end{center}

\subsection{Simulation Interaction Standards}

Every simulation across all acts must follow a consistent structure:
\begin{enumerate}[nosep]
  \item \textbf{What is being simulated} --- one-sentence description.
  \item \textbf{Controls} --- sliders, toggles, and buttons.
  \item \textbf{Visual outputs} --- panels, charts, counters.
  \item \textbf{Key takeaway} --- what the user should learn.
\end{enumerate}

All simulations must support:
\begin{itemize}[nosep]
  \item \textbf{Interactive controls:} sliders and toggles respond in
        real time.
  \item \textbf{Fixed seed toggle:} deterministic replay for reproducible
        exploration.
  \item \textbf{``Run 1,000 experiments'' mode:} batch mode showing
        aggregate statistics (false positive rate, average stopping time,
        power).
  \item \textbf{``Reset simulation'' button:} returns all parameters and
        visuals to their default state.
\end{itemize}

\subsection{Tooltips}

Hoverable tooltips appear on key quantities throughout all acts:
\begin{itemize}[nosep]
  \item \textbf{CI:} ``A range of plausible values for the treatment effect,
        valid at every moment since the experiment started.''
  \item \textbf{Lift ($\hat{\tau}(t) / \hat{\mu}_0(t)$):} ``The percentage
        change in the metric caused by the treatment, relative to control.''
  \item \textbf{Variance:} ``How spread out individual outcomes are.  Higher
        variance means more noise and slower experiments.''
  \item \textbf{Sequential multiplier:} ``The factor that replaces 1.96 in a
        classical CI to guarantee validity under continuous monitoring.''
\end{itemize}

\subsection{Transitions Between Acts}

Acts transition with smooth scroll animations:
\begin{itemize}[nosep]
  \item Clicking ``Next act'' smoothly scrolls to the next section.
  \item The persistent sidebar highlights the current act.
  \item Simulation state from the previous act can optionally carry forward
        (e.g.\ the same dataset used in Acts~10--12).
\end{itemize}


% ══════════════════════════════════════════════════════════════════════════════
\newpage
\section{Summary: The Through-Line}

\begin{center}
\renewcommand{\arraystretch}{1.3}
\begin{tabular}{@{} c l l p{6.5cm} @{}}
\toprule
\textbf{Act} & \textbf{Object} & \textbf{Represents} & \textbf{Key idea} \\
\midrule
0 & Repeated $p$-values
  & The peeking problem
  & Checking repeatedly inflates FP \\
1 & Random walk $S_n$
  & Position
  & $\SD$ grows as $\sqrt{n}$ \\
2 & Winnings $M_n$
  & Gambler's fortune
  & Martingale; Doob: $\E[M_\tau]=M_0$ \\
3 & Likelihood ratio $\LR_n$
  & Evidence against $H_0$
  & Comparing hypotheses flip by flip \\
4 & $\LR_n$ under $H_0$
  & LR is a martingale
  & The cancellation proof; $\LR_0{=}1$, $\LR_n{\geq}0$ \\
5 & Markov $+$ Ville
  & Anytime-valid guarantee
  & $\Prob(\LR_n \text{ ever}{\geq}1/\alpha)\leq\alpha$ \\
6 & SPRT
  & Wald's decision rule
  & Optimal but requires known $\delta$ \\
7 & $\LR_n^{H}$ (mixture LR)
  & Robust evidence
  & Robbins' mixtures $\to$ Johari's mSPRT \\
8 & Confidence sequence
  & Modern framework
  & Howard: nonparametric, estimated $\sigma^2$ \\
9 & Peeking vs.\ variance
  & Three problems
  & Peeking, high variance, heterogeneous users \\
10 & Eppo pipeline
  & Step-by-step factory
  & Randomise $\to$ adjust $\to$ CI $\to$ decide \\
11 & Regression adjustment
  & Noise removal
  & Generalised CUPED; $\Var(Y^*)=\Var(Y)(1-R^2)$ \\
12 & Eppo's sequential CI
  & Final output
  & Sequential multiplier; $\nu=M\hat{\sigma}^2$; decide anytime \\
13 & DIY alternatives
  & Simple approximations
  & Bonferroni, Pocock, OBF vs.\ Eppo \\
\bottomrule
\end{tabular}
\end{center}

\bigskip
\textbf{The core message:} All fourteen acts build on the same mathematical
structure --- non-negative (super)martingales and Ville's inequality.  The
anytime-valid guarantee flows from a single elegant fact:
\[
  \textit{A fair game cannot be beaten by choosing when to quit.}
\]

The evolution across eight decades:
\[
  \underset{1939}{\text{Ville}}
  \;\to\;
  \underset{1945}{\text{Wald (SPRT)}}
  \;\to\;
  \underset{1970}{\text{Robbins (mixtures)}}
  \;\to\;
  \underset{2017}{\text{Johari (mSPRT)}}
  \;\to\;
  \underset{2021}{\text{Howard (CS)}}
  \;\to\;
  \underset{2024}{\text{Eppo}}
\]


% ══════════════════════════════════════════════════════════════════════════════
\newpage
\section{References}
\label{sec:references}

\begin{enumerate}
  \item \textbf{Ville, J.} (1939).
        \textit{\'Etude critique de la notion de collectif}.
        Gauthier-Villars, Paris.
        --- Original proof of Ville's inequality for non-negative
        supermartingales.
        \textit{[Used in Acts~5, 7, 8.]}

  \item \textbf{Wald, A.} (1945).
        Sequential tests of statistical hypotheses.
        \textit{Annals of Mathematical Statistics}, 16(2), 117--186.
        \url{https://doi.org/10.1214/aoms/1177731118}
        --- Foundational 70-page monograph defining the SPRT, proving
        the fundamental inequalities for error probabilities, deriving the
        ASN formulas, and showing 47--63\% savings over fixed-sample tests.
        \textit{[Used in Act~6.]}

  \item \textbf{Wald, A.\ \& Wolfowitz, J.} (1948).
        Optimum character of the sequential probability ratio test.
        \textit{Annals of Mathematical Statistics}, 19(3), 326--339.
        \url{https://doi.org/10.1214/aoms/1177730197}
        --- Proof that the SPRT is exactly optimal: among all sequential tests
        with the same error probabilities, it minimises expected sample size
        under both $H_0$ and $H_1$.
        \textit{[Used in Act~6.]}

  \item \textbf{Doob, J.\,L.} (1953).
        \textit{Stochastic Processes}.
        Wiley, New York.
        --- Definitive treatment of martingale theory including the
        Optional Stopping Theorem.
        \textit{[Used in Act~2.]}

  \item \textbf{Robbins, H.} (1970).
        Statistical methods related to the law of the iterated logarithm.
        \textit{Annals of Mathematical Statistics}, 41(5), 1397--1409.
        \url{https://doi.org/10.1214/aoms/1177696786}
        --- Introduced the mixture likelihood ratio approach and confidence
        sequences.  Showed that weighting over alternatives via a mixing
        distribution produces a martingale amenable to Ville's inequality.
        Derived closed-form boundaries for the Normal case (including the root
        of the Normal mixture that Howard et al.\ later generalised) and
        connected boundary growth rates to the law of the iterated logarithm.
        \textit{[Used in Act~7.]}

  \item \textbf{Deng, A., Xu, Y., Kohavi, R., \& Walker, T.} (2013).
        Improving the sensitivity of online controlled experiments by
        utilizing pre-experiment data.
        \textit{Proc.\ WSDM}, 123--132.
        \url{https://doi.org/10.1145/2433396.2433413}
        --- The original CUPED paper.  Optimal adjustment coefficient is the
        OLS slope $\theta^* = \Cov(Y,X)/\Var(X)$; variance is reduced by a
        factor of $1-\rho^2$.  Practical results at Microsoft Bing: 45--52\%
        variance reduction with 2-week pre-period.
        \textit{[Used in Acts~9--11.]}

  \item \textbf{Johari, R., Pekelis, L., \& Walsh, D.J.} (2017).
        Always valid inference: Continuous monitoring of A/B tests.
        \textit{Operations Research}, 65(6), 1651--1667.
        \url{https://doi.org/10.1287/opre.2017.1607}
        --- Applied the mixture SPRT to A/B testing.  Derived the closed-form
        mSPRT for Normal observations, always-valid $p$-values and confidence
        intervals, and guidance for calibrating $\tau$.  Assumes known
        variance.
        \textit{[Used in Act~7.]}

  \item \textbf{Howard, S.\,R., Ramdas, A., McAuliffe, J., \& Sekhon, J.}
        (2021).
        Time-uniform, nonparametric, nonasymptotic confidence sequences.
        \textit{Annals of Statistics}, 49(2), 1055--1080.
        \url{https://doi.org/10.1214/20-AOS1991}
        --- Developed the sub-$\psi$ framework for confidence sequences.
        Key result: the Normal mixture boundary
        $u(v) = \sqrt{(v+\nu)\log\frac{v+\nu}{\nu\alpha^2}}$.
        Nonparametric (works beyond Normal data), nonasymptotic (finite-sample
        guarantees), and allows estimated variance.  \textbf{This is the
        framework Eppo adopted.}
        \textit{[Used in Acts~8, 10, 12.]}

  \item \textbf{Schmit, M.\ \& Miller, T.} (2024).
        Sequential confidence intervals for comparing two proportions with
        variance reduction.
        \textit{Eppo Technical Report}.
        --- Eppo's implementation.  Combines Howard et al.'s Normal mixture
        confidence sequence with regression-adjusted variance reduction
        (generalised CUPED).  Estimates variance (does not assume known
        $\sigma^2$).  Tuning parameter is $\nu$ (set from expected sample
        size $M$).  Reports relative lift with sequential confidence intervals.
        \textit{[Used in Acts~9--12.]}
\end{enumerate}


% ══════════════════════════════════════════════════════════════════════════════
\newpage
\section*{Appendix: What Changed from v3}
\addcontentsline{toc}{section}{Appendix: What Changed from v3}

\begin{enumerate}[nosep]
  \item \textbf{Major correction:} Eppo does \emph{not} use the mSPRT
        (Johari et al.\ 2017).  They use the confidence sequence framework
        of Howard et al.\ (2021), specifically the Normal mixture boundary.
        Acts~8 and~9 completely rewritten.
  \item \textbf{GAVI removed.}  The term ``GAVI'' does not appear in the
        Schmit \& Miller paper.  The tuning parameter is $\nu$, set from the
        expected sample size $M$, not from an empirical distribution of past
        effect sizes.
  \item \textbf{Robbins (1970) properly credited.}  His mixture approach and
        confidence sequences are now covered in Act~7A, before Johari's mSPRT
        (Act~7B).  This corrects the historical attribution.
  \item \textbf{Act~6 enriched:} Added Wald's ASN formula, savings table
        (47--63\%), the random walk view, truncation, and the optimality
        conjecture resolved by Wald \& Wolfowitz (1948).
  \item \textbf{CUPED corrected:} Added the validity condition
        ($\E[X^{(t)}] = \E[X^{(c)}]$), the warning against post-treatment
        covariates (from Deng et al.\ 2013), recommended pre-period length,
        practical results from Bing.
  \item \textbf{Notation cheat sheet removed:} All notation is now explained
        inline at first use.
  \item \textbf{Act numbering:}  Now 14 acts (0--13) instead of 10 (0--9).
        Former Act~9 (Eppo, all-in-one) expanded into Acts~9--12 following
        detailed guidance: Act~9 (problem Eppo solves), Act~10 (pipeline
        step-by-step), Act~11 (variance reduction via regression), Act~12
        (sequential CIs in Eppo).  Act~13 added: DIY alternatives
        (Bonferroni, Pocock, O'Brien--Fleming) with head-to-head Eppo
        comparison.
  \item \textbf{References:}  Each reference now notes which acts use it and
        describes the specific contribution more precisely.
  \item \textbf{Historical notes} added throughout (Ville's 1939 thesis,
        Wald's 1945 optimality conjecture, Wald--Wolfowitz 1948 resolution).
  \item \textbf{Through-line updated:} Timeline from Ville (1939) through
        Eppo (2024) added to the summary.
\end{enumerate}

\end{document}
